--- Page 1 ---
 
 
S1 
 
 
Supporting Information 
 
Low-spin 1,1ʹ-diphosphametallocenates of Chromium and Iron 
Samuel M. Greer,1,2,3 Ökten Üngör,3 Ross J. Beattie,1 Jaqueline L. Kiplinger,*1 Brian L. Scott,1 Benjamin W. Stein,*1 and 
Conrad A. P. Goodwin*1,4 
 
1Chemistry Division, Los Alamos National Laboratory, Los Alamos, NM 87545, USA. 
2National High Magnetic Field Laboratory, Florida State University, Tallahassee, FL 32310, USA. 
3Department of Chemistry & Biochemistry, Florida State University, Tallahassee, FL 32306, USA. 
4Department of Chemistry, School of Natural Sciences, The University of Manchester, Oxford Road, Manchester, M13 9PL, 
UK. 
 
*Correspondence to: kiplinger@lanl.gov; bstein@lanl.gov; cgoodwin@lanl.gov.  
 
 
Electronic Supplementary Material (ESI) for ChemComm.
This journal is © The Royal Society of Chemistry 2020


--- Page 2 ---
 
 
S2 
 
 
Table of Contents 
S1. Experimental Details 
 
 
 
 
 
 
 
 
 
 
 
 
3 
 
S1.1.  
Preparation of KTMP 
 
 
 
 
 
 
 
 
 
 
4 
 
S1.2.  
Synthesis of 2-M (M = Cr, Fe)  
 
 
 
 
 
 
 
 
5 
 
S1.3.  
Attempted extension to Mn and Co  
 
 
 
 
 
 
 
6 
S2. Crystallography 
 
 
 
 
 
 
 
 
 
 
 
 
 
8 
S3. Molecular Structures 
 
 
 
 
 
 
 
 
 
 
 
 
11 
S4. NMR Spectroscopy  
 
 
 
 
 
 
 
 
 
 
 
 
17 
S5. UV-vis-nIR Spectroscopy  
 
 
 
 
 
 
 
 
 
 
 
23 
S6. ATR-FTIR Spectroscopy  
 
 
 
 
 
 
 
 
 
 
 
27 
S7. EPR Spectroscopy  
 
 
 
 
 
 
 
 
 
 
 
 
29 
S8. 57Fe Mössbauer Spectroscopy  
 
 
 
 
 
 
 
 
 
 
31 
S9. Quantum Chemical Calculations 
 
 
 
 
 
 
 
 
 
 
32 
 
S9.1.  
Example ORCA input files 
 
 
 
 
 
 
 
 
 
33 
S10. SQUID Magnetometry  
 
 
 
 
 
 
 
 
 
 
 
35 
S11. Qualitative ligand HOMOs in Cp– vs {PC4H4}– 
 
 
 
 
 
 
 
36 
S12. References 
 
 
 
 
 
 
 
 
 
 
 
 
 
 
37 
S13. CIF reports 
 
 
 
 
 
 
 
 
 
 
 
 
 
 
39 
 
 


--- Page 3 ---
 
 
S3 
 
 
S1. Experimental Details 
General considerations 
All syntheses and manipulations were conducted under UHP argon (AirGas) with rigorous exclusion of oxygen and water 
using Schlenk line and glove box techniques. Molecular sieves (4 Å, Sigma Aldrich, 8-12 mesh) were activated at 200 °C for 
12 hours, at 5x10–4 mbar. All solvents were degassed (by bubbling Ar through them, or under dynamic vacuum) prior to drying 
procedures. Anhydrous THF (Sigma Aldrich) was degassed, refluxed over K/Ph2CO until a persistent purple color was 
obtained, then distilled and stored over 4 Å molecular sieves. Anhydrous hexane and toluene (Sigma Aldrich) were degassed, 
refluxed over NaK/Ph2CO for several days until dark coloration (purple solutions or purple/green solids) were observed, then 
distilled and stored over 4 Å molecular sieves. All solvents were tested with a dilute THF solution of Na2Ph2CO (150 mg 
Ph2CO in 20 mL of THF with an excess of Na metal) such that ethereal solvents required 1 drop / mL to retain purple coloration 
and hydrocarbon solvents required 1 drop / 2 mL. Anhydrous C6D6 and C4D8O (both Sigma Aldrich) for NMR spectroscopy 
were degassed, stirred with NaK for 72 hours and filtered through a glass-fiber filter disc prior to use. KC8 was prepared by 
heating graphite flakes (200 mesh, Sigma Aldrich) and 1/8th of a molar equivalent of freshly cleaned K metal together in a 
glass scintillation vial inside an argon glovebox with vigorous stirring using a metal spatula. 2-butyne was vacuum transferred 
and stored over 4 Å sieves for 4 days prior to use. 2.2.2-cryptand (Alfa Aesar) was dried under vacuum (5x10–4 mbar) for 48 
hours and stored in a glovebox. PhPCl2, anhydrous CrCl2 (Alfa Aesar), anhydrous FeCl2 (Alfa Aesar) and anhydrous CoCl2 
(Fisher Scientific) were used as received. [M(TMP)2] (1-M, M = Cr or Fe) were prepared using a procedure slightly modified 
from those previously described;1,2 anhydrous FeCl2 and CrCl2 were combined with base free KTMP in THF at room 
temperature, and the complexes were crystallized from hexane. The glovebox atmosphere was periodically checked using a 
dilute toluene solution of [Ti(Cp)2(μ-Cl)]2 (200 mg of commercial [Ti(Cp)2(Cl)2] reduced over an excess of Zn powder in 20 
mL of toluene, and filtered. A drop of the green solution was allowed to dry on a vial lid (Urea, VWR) prior to any 
manipulations, such that the residue dried to a dark green color each time (a color change to yellow or orange indicates 
decomposition of the Ti test compound and atmospheric O2/H2O removal is required). All glassware, and glass-fiber filter 
discs, were stored in an oven (>150 °C) for 24 hours prior to being brought into the glovebox, and FEP (fluorinated ethylene 
propylene) NMR liners were brought into the box via overnight vacuum cycles. The NMR solutions were loaded into a fresh 
FEP NMR liner, stoppered with a PTFE plug, and then placed in J. Young tap appended NMR tubes. This was done as it 
facilitates the easy handling of material in facilities primarily concerned with Th/dU work and is a local requirement for the 
area this work had to be performed in. Crystals for single-crystal X-ray diffraction were mounted on nylon loops in Paratone-
N oil. Solution phase electronic absorption spectra were collected at ambient temperature using a Varian Cary 6000i UV-vis-


--- Page 4 ---
 
 
S4 
 
 
nIR spectrometer. Solutions were contained in low volume (1 mL) screw-capped quartz cuvettes (1 cm path length). ATR-
FTIR spectra were obtained using a Thermo Scientific Nicolet iS5 FTIR spectrometer using a Golden Gate Diamond ATR 
(ZnSe lenses) with a reaction anvil. Experimental details for other techniques (57Fe Mössbauer and EPR spectroscopies, and 
SQUID magnetometry) are described separately in their corresponding sections. 
 
S1.1. Preparation of KTMP 
Base-free KTMP was prepared by a multi-step procedure adapted from several literature procedures and beginning with cheap, 
commercially available starting materials (Scheme 1). Here we have structurally characterized two of the intermediates, 
[Zr(Cp)2(C4Me4) (5) and Ph-PC4Me4 (Ph-TMP, 6) for the first time. 
 
 
Scheme S1. Synthesis of KTMP, beginning from dicyclopentadiene / Na and anhydrous ZrCl4. 
 
[Zr(Cp)2(Cl)2] was prepared from [Zr(Cl)4(THF)2] (prepared from ZrCl4 and stoichiometric THF in DCM, and recrystallized 
from DCM)3 and two equivalents of NaCp (prepared from dicyclopentadiene and Na metal) in THF,4 followed by solvent 
removal and extraction of the bright yellow residue with DCM. [Zr(Cp)2(C4Me4)] (5) was prepared from [Zr(Cp)2(Cl)2], nBuLi, 
and 2-butyne, and crystallized from hexane or used in situ.5,6 Ph-TMP (6) was prepared from 5 and PhPCl2, then purified by 
distillation which afforded colorless crystals, rather than a colorless oil.5 During the course of this study we have structurally 
characterized 5 and 6 for the first time (vide infra). Base-free KTMP was prepared from 6 and a four-fold excess of K metal 
in boiling DME or THF, followed by filtration of the bright orange solution and exhaustive extraction of the gray/brown 
ZrCl4 + 2 THF
DCM, 0°C
Excess CpH + Na
Neat, �
NaCp
[Zr(Cl)4(THF)2]
2 x
THF, –78°C
[Zr(Cp)2(Cl)2]
[Zr(Cp)2(Cl)2] + 2.3 MeCCMe
THF, –78°C
[Zr(Cp)2(C4Me4)2]
2.1 nBuLi
THF, –78°C
2.2 PhPCl2 Ph-PC4Me4
5
6
+ [Zr(Cp)2(Cl)2]
This procedure all performed in one flask with warming and re-cooling steps.
6 + 4 K
THF or DME, �
Solvated KTMP + ~ 2 K + “KPh” decomposition products


--- Page 5 ---
 
 
S5 
 
 
residue. Multiple crops of colorless crystals were collected by repeated concentration and cooling of the orange solution. This 
material was then fully desolvated in vacuo to afford a white powder, which was confirmed to be base-free by elemental 
analysis. (TMP)2 was prepared as previously described, and purified by crystallization from hexane.5 
 
S1.2. Synthesis of 2-M (M = Cr, Fe) 
Synthesis of [K(2.2.2-crypt)][Cr(TMP)2] 2-Cr: A solution of 1-Cr (55 mg, 0.17 mmol) and 2.2.2-crypt (62 mg, 0.17 mmol) 
in THF (1 mL) was prepared in a 20 mL glass scintillation vial equipped with a Teflon stir bar, then cooled to –35°C for 1 
hour. Separately, a vial of KC8 was also chilled to –35°C along with an aluminum vial block. Cold KC8 (25 mg, 0.18 mmol, 
1.1 eq.) was added rapidly to the vigorously stirred THF solution which was maintained close to –35°C by the vial block. The 
solution immediately turned from dark red to dark brown/black. The brown/black solution was filtered from black solids 
(presumably graphite) through a glass pipette loaded with a glass-fiber disc padded with a 2 mm layer of KC8 into a 4 mL 
glass vial. Hexane (3 mL) was carefully layered on top of the THF solution, and the vial was then stored at –35°C for 18 hours. 
The pale brown/colorless supernatant was decanted and the brown/black crystals were washed with room temperature hexane 
(3 x 2 mL) followed by drying in vacuo to afford 2-Cr (yield 102 mg, 82%). CrKP2N2O6C34H60; calcd. C 54.75, H 8.11, N 
3.76; found C 54.32, H 8.18, N 3.49. 1H NMR (C4D8O, 400 MHz, 298 K): 2.63, 3.62, 3.65, only peaks attributable to free 
2.2.2-cryptand, and to the solvent residual signals were observed. 31P{1H} NMR (C4D8O, 162 MHz, 298 K): 74.55, 
diamagnetic impurity, corresponds to free {TMP}– anion. ATR-IR (𝜈̅, cm–1): 2958 (m), 2935 (w), 2922 (w), 2866 (m), 2825 
(s), 2727 (w), 2698 (w), 2360 (m), 2343 (m), 1477 (m), 1454 (w), 1442 (m), 1412 (w), 1398 (w), 1369 (m), 1358 (m), 1350 
(s), 1300 (m), 1292 (m), 1257 (m), 1240 (w), 1169 (w), 1132 (s), 1101 (vs), 1080 (vs), 1057 (m), 1018 (m), 947 (s), 931 (s), 
831 (m), 820 (m), 752 (m), 681 (w), 669 (w), 658 (vw), 571 (m). UV-vis-nIR (THF) λmax (cm–1, ε): 502 (19904, 4180, vbr). 
 
Synthesis of [K(2.2.2-crypt)][Fe(TMP)2] 2-Fe: A solution of 1-Fe (150 mg, 0.45 mmol) and 2.2.2-crypt (170 mg, 0.45 mmol) 
in THF (1 mL) was prepared in a 20 mL glass scintillation vial equipped with a Teflon stir bar, then cooled to –35°C for 1 
hour. Separately, a vial of KC8 was also chilled to –35°C along with an aluminum vial block. Cold KC8 (67 mg, 0.50 mmol, 
1.1 eq.) was added rapidly to the vigorously stirred THF solution which was maintained close to –35°C by the vial block. The 
solution immediately turned from dark red to dark green/black. The green/black solution was filtered from black solids 
(presumably graphite) through a glass pipette loaded with a glass-fiber disc padded with a 2 mm layer of KC8 into a 4 mL 
glass vial. Hexane (3 mL) was carefully layered on top of the THF solution, and the vial was then stored at –35°C for 18 hours. 
The pale brown/colorless supernatant was decanted and the green/black crystals were washed with room temperature hexane 


--- Page 6 ---
 
 
S6 
 
 
(3 x 2 mL) followed by drying in vacuo to afford 2-Fe (yield 243 mg, 72%). FeKP2N2O6C34H60; calcd. C 54.47, H 8.07, N 
3.74; found C 54.96, H 8.35, N 3.61. 1H NMR (C4D8O, 400 MHz, 298 K): 1.94 (s, 12H, TMP-3,4-Me2), 2.19 (d, 3JHP = 10.29 
Hz, 12H, TMP-2,5-Me2), 2.28 (s, 12H, 2.2.2-crypt), 3.30 (s, 12H, 2.2.2-crypt), 3.37 (s, 24H, 2.2.2-crypt). 31P{1H} NMR 
(C4D8O, 162 MHz, 298 K): 74.46, diamagnetic impurity, corresponds to free {TMP}– anion. ATR-IR (𝜈̅, cm–1): 2954 (w), 
2873 (m), 2835 (m), 2729 (w), 2698 (w), 2360 (w), 2343 (w), 2324 (w), 1475 (m), 1442 (m), 1414 (w), 1390 (w), 1358 (m), 
1350 (m), 1296 (m), 1257 (m), 1240 (w), 1173 (w), 1132 (s), 1099 (vs), 1078 (vs), 1020 (m), 947 (s), 931 (s), 829 (m), 820 
(m), 796 (w), 752 (m), 679 (w), 667 (w), 638 (vw), 625 (vw), 588 (w), 569 (w), 553 (w), 536 (w). UV-vis-nIR (THF) λmax 
(cm–1, ε): 428 (23375, 2220, vbr), 533 (18754, 1830, vbr), 616 (16239, 1720, vbr), 885 (11297, 1760), 924 (10818, 1730). 
 
S1.3. Attempted extension to Mn and Co 
The synthesis of a redox series, such as that for 3-M and 4-M (M = Mn, Fe, Co) was desirable. Attempts to synthesize the 
required 1-Mn from commercial MnCl2, MnCl2(THF)n prepared in-situ, MnCl2(DME)n prepared in-situ, or MnCl2(PMe3) with 
two equivalents of KTMP in THF, DME or toluene resulted in intractable mixtures of dark powder that were insoluble in all 
common solvents that are compatible with this ligand set. The synthesis of 1-Co was attempted using anhydrous CoCl2 and 
two equivalents of K-TMP in THF (see below). 
 
Synthesis of [Co(η5-TMP)(μ2:η1:η1-TMP)]2 7: CoCl2 (65 mg, 0.5 mmol) was slurried in THF (4 mL) in a 20 mL glass 
scintillation vial at room temperature. K-TMP (178 mg, 1 mmol) was added to the blue suspension, which immediately turned 
dark brown. The brown mixture was stirred for 16 hours at room temperature. Volatiles were removed in vacuo and the black 
tacky residue was extracted with hot (60 °C) hexane (6 mL) which afforded a dark brown/black solution which was 
concentrated to 1.5 mL and stored at room temperature for 16 hours to afford black crystals of 7 along with sticky brown 
residues (mass 14 mg). We were unable to optimize the yield of this complex. A 31P{1H} NMR taken of the crystalline material 
indicated several products (Figure S21). 
 
Attempted reduction of crude 7: The reaction above was repeated on the same scale. We then attempted reduction of the 
crude mixture, assuming an arbitrary 80% conversion of the reactants to 7, presuming any reduced product might form a salt-
like complex as 2-M, and be separable from 7 based on solubility. After stirring for 16 hours at room temperature, the crude 
mixture was reduced to dryness in vacuo, and extracted with warm (45 °C) toluene (5 mL), filtered through a glass pipette 
loaded with a glass-fiber disc into a 20 mL glass scintillation vial, and the brown/black solution was reduced to dryness in 


--- Page 7 ---
 
 
S7 
 
 
vacuo. Solid 2.2.2-crypt (151 mg, 0.4 mmol, 0.8 eq. assuming 80% conversion) was added along with a Teflon stir bar, and 
then the tacky black residue was dissolved in THF (2 mL), and cooled to –35 °C. Pre-cooled (–35 °C) solid KC8 (54 mg, 0.40 
mmol, 0.8 eq. assuming 80% conversion) was added to the rapidly stirred solution which was maintained at –35°C using a 
chilled vial block. No perceptible color change was observed. After stirring for 1 minute the black solution was filtered 
underneath pre-cooled hexane (10 mL) in a 20 mL glass scintillation vial, and stored overnight at –35°C which afforded several 
low quality crystals of [K(2.2.2-crypt)][TMP] (8), 2.2.2-crypt, and no other crystalline products. The brown solution was 
decanted from the crystals and reduced to a thin oil in vacuo, which formed biphasic mixtures in hydrocarbon solvents (pentane, 
hexane, toluene), and pale brown solutions in ethereal solvents (THF, DME, Et2O), but no further crystalline products could 
be isolated. When the reaction was repeated at room temperature, the observed colors and biphasic mixtures etc. were all 
identical, and likewise no crystalline products that contained Co could be isolated in our hands. 
 
 


--- Page 8 ---
 
 
S8 
 
 
S2. Crystallography 
General considerations 
The crystal data for complexes 1-Cr, 1-Fe, 2-Cr, 2-Fe, 5 – 8 are compiled in Tables S1-3. Crystals of 1-Fe, 2-Cr, 2-Fe, 7, and 
8 were examined with a Bruker D8 Quest diffractometer equipped with a CMOS detector and using mirror-monochromated 
Mo Kα radiation (λ = 0.71073 Å) operating in shutterless mode; crystals of 1-Cr was examined with a Bruker APEX II 
diffractometer equipped with a CCD detector and using mirror-monochromated Mo Kα radiation (λ = 0.71073 Å); crystals of 
5 were examined using an Oxford Diffraction Supernova diffractometer, equipped with CCD area detector and a mirror-
monochromated Mo Kα radiation (λ = 0.71073 Å); and crystals of 4 were examined with a Rigaku XtalLAB AFC11 
diffractometer, equipped with CCD detector and mirror-monochromated Cu Kα radiation (λ = 1.54184 Å). APEX II (APEX 
II), or APEX III (D8 Quest), or CrysAlisPro (Rigaku XtalLAB AFC11) software was used for control and solving the unit 
cells prior to data collection. Intensities were integrated from data recorded on 0.5° frames by ω rotation with 10s frame 
exposure (1-Cr), 13s frame exposure (1-Fe), 25s frame exposure (2-Fe), 30s frame exposure (2-Cr), or 90s frame exposure; 
or by 0.7° frames by ω rotation with 1s or 2s frame exposure (6); or finally by 0.8° frames by ω rotation with 30s frame 
exposure (5). CrysAlisPro was used for final unit cell determination and parameters were refined from the observed positions 
of all strong reflections in each data set and an analytical absorption correction was applied.7 The Olex28 GUI was used for 
structure solution and refinement utilizing the ShelX software packages.9,10 The structures were solved using ShelXT;10 the 
datasets were refined by ShelXL9 using full-matrix least-squares on all unique F2 values, with anisotropic displacement 
parameters for all non-hydrogen atoms, and with constrained riding hydrogen geometries; Uiso(H) was set at 1.2 (1.5 for methyl 
groups) times Ueq of the parent atom. The largest features in final difference syntheses were close to heavy atoms and were 
of no chemical significance. Olex2 combined with Inkscape was employed for molecular graphics.8,11 CCDC 2003385 (1-Cr), 
2003386 (1-Fe), 2003387 (2-Cr), 2003388 (2-Fe), 2003389 (5), 2003390 (6), 2017380 (7), and 2017381 (8) contain the 
supplementary crystal data for this article.§ These data can be obtained free of charge from the Cambridge Crystallographic 
Data Centre via www.ccdc.cam.ac.uk/data_request/cif. 
 
 
 
 
 
§ LA-UR-20-20981 


--- Page 9 ---
 
 
S9 
 
 
Table S1. Crystallographic data for 1-Cr, 1-Fe, 2-Cr, and 2-Fe. 
 
1-Cr 
1-Fe 
2-Cr 
2-Fe 
 
Internal identifier 
Apx2920 
Dq0894 
Dq1027 
Dq0907 
 
CCDC ref code 
2003385 
2003386 
2003387 
2003389 
 
Formula 
CrP2C16H24 
FeP2C16H24 
KCrP2N2O6C34H60 
KFeP2N2O6C34H60 
 
Fw 
330.29 
334.14 
745.88 
749.73 
 
Crystal syst 
Monoclinic 
Monoclinic 
Triclinic 
Triclinic 
 
Space group 
P21/n 
C2/c 
P–1 
P–1 
 
a, Å 
7.8420(3) 
14.2948(7) 
10.1994(4) 
10.1947(2) 
 
b, Å 
12.5620(4) 
12.8549(6) 
14.2193(7) 
14.2373(4) 
 
c, Å 
8.8611(3) 
8.8448(4) 
15.0882(6) 
15.0677(3) 
 
α, ° 
90 
90 
62.328(4) 
62.275(2) 
 
β, ° 
109.444(4) 
104.364(5) 
86.597(3) 
85.840(2) 
 
γ, ° 
90 
90 
87.156(3) 
87.045(2) 
 
V, Å3 
823.13(5) 
1574.50(13) 
1933.91(16 
1930.42(8) 
 
Z 
2 
4 
2 
2 
 
ρcalcd, g cm–3 
1.333 
1.410 
1.281 
1.290 
 
μ, mm-1 
0.874 
1.146 
0.529 
0.624 
 
F(000) 
348 
704 
798 
802 
 
Cryst size, mm 
0.30 x 0.35 x 0.42 
0.16 x 0.19 x 0.29 
0.30 x 0.30 x 0.30 
0.05 x 0.22 x 0.25 
 
Temperature, K 
100(2) 
100(2) 
100(2) 
100(2) 
 
no. reflections (unique) 
6774 (1508) 
6215 (1605) 
23292 (7060) 
25573 (7807) 
 
Rint 
0.017 
0.032 
0.033 
0.036 
 
R1(wR2) (F2 > 2σ(F2)) 
0.0252 (0.0667) 
0.0312 (0.0789) 
0.0338 (0.0872) 
0.0400 (0.1035) 
 
Sa 
1.09 
1.12 
1.10 
1.09 
 
min./max. diff map, Å–3 
–0.26, 0.33 
–0.24, 0.56 
–0.25, 0.49 
–0.26, 0.62 
 
a R = ∑||FO| – |FC||/∑|FO|; RW = [∑w(FO2 – FC2)2/∑w(FO2)2]0.5; S = [∑w(FO2 – FC2)2/(no. data – no. params)]0.5 for all data. 
 
 


--- Page 10 ---
 
 
S10 
 
 
Table S2. Crystallographic data for 5 – 8. 
 
5 
6 
7 
8 
Internal identifier 
Adpm457 
Ldpm35 
Dq1025 
Dq0919 
Formula 
ZrC18H22 
PC14H17 
Co2P4C32H48 
K2P2O12N4C52H96·C4H8O 
Fw 
329.57 
216.24 
674.44 
1181.57 
Crystal syst 
Monoclinic 
Triclinic 
Monoclinic 
Triclinic 
Space group 
I2/a 
P–1 
C2/c 
P–1 
a, Å 
14.1554(8)  
8.0727(2) 
20.161(3) 
12.7189(8) 
b, Å 
9.6976(4) 
8.74422(19) 
9.0620(4) 
15.9437(10) 
c, Å 
12.0182(6) 
18.4065(4) 
20.214(3) 
17.8937(12) 
α, ° 
90 
103.1739(17) 
90 
64.066(6) 
β, ° 
112.336(6) 
92.6248(19) 
119.71(2) 
85.655(5) 
γ, ° 
90 
100.705(2) 
90 
86.321(5) 
V, Å3 
1526.00(15) 
1237.81(5) 
3207.6(9) 
3251.9(4) 
Z 
4 
4 
4 
2 
ρcalcd, g cm–3 
1.434 
1.160 
1.397 
1.207 
μ, mm-1 
0.705 
1.664 
1.254 
0.254 
F(000) 
680 
464 
1416 
1280 
Cryst size, mm 
0.04 x 0.12 x 0.17 
0.09 x 0.18 x 0.19 
0.10 x 0.12 x 0.40 
0.05 x 0.14 x 0.37 
Temperature, K 
150(2) 
150(2) 
100(2) 
100(2) 
no. reflections (unique) 
2872 (1558) 
13067 (4482) 
17531 (2935) 
31744 (13040) 
Rint 
0.028 
0.019 
0.065 
0.065 
R1(wR2) (F2 > 2σ(F2)) 
0.0330 (0.0805) 
0.0329 (0.0899) 
0.0333 (0.0834) 
0.0673 (0.1979) 
Sa 
1.07 
1.04 
1.01 
1.06 
min./max. diff map, Å–3 
–0.54, 0.46 
–0.26, 0.32 
–0.34, 0.50 
–0.54, 0.73 
a R = ∑||FO| – |FC||/∑|FO|; RW = [∑w(FO2 – FC2)2/∑w(FO2)2]0.5; S = [∑w(FO2 – FC2)2/(no. data – no. params)]0.5 for all data. 
 
 


--- Page 11 ---
 
 
S11 
 
 
S3. Molecular structures 
 
Figure S1. Side view (left) and top view (right) of the molecular structure of 1-Cr with ellipsoids set at 50%, and hydrogen 
atoms removed for clarity. Cr(1) sits on a site of symmetry, and thus the inter-ring angle is necessarily 180°. (operations: x, y, 
z; 1–x, 1–y, 1–z). 
 
 
Figure S2. Side view (left) and top view (right) of the molecular structure of 1-Fe with ellipsoids set at 50%, and hydrogen 
atoms removed for clarity. Fe(1) sits on a site of symmetry, and thus the inter-ring angle is necessarily 180°. (operations: x, y, 
z; ½–x, ½–y, 1–z). 
Table S3. Structural parameters for 1-Cr and 1-Fe. 
M = Cr or Fea 
1-Cr 
1-Fe 
P–M / Å 
2.3812(4) 
2.2932(4) 
TMPcent···M / Å 
1.795(1) 
1.660(1) 
PC2plane···C4planeb / ° 
4.45(12) 
1.27(13) 
a Prior-reported CCDC codes for 1-Cr: NORRAN; 1-Fe: ABIGAU. 
b A plane is defined by the three atoms, P and the two adjoining C atoms; a second plane is defined by all four C atoms of the 
TMP ring. The “hinge angle”, HA, between these two planes is reported. 
 
Cr(1)
P(1)
P(1A)
Cr(1)
P(1)
P(1A)
Fe(1)
P(1)
P(1A)
Fe(1)
P(1)
P(1A)


--- Page 12 ---
 
 
S12 
 
 
 
Figure S3. Molecular structure of 2-Cr with ellipsoids set at 50%, and hydrogen atoms removed for clarity. The structure has 
been inverted for this image to show the similarity to 2-Fe (below). 
 
 
Figure S4. Molecular structure of 2-Fe with ellipsoids set at 50%, and hydrogen atoms removed for clarity. 
 
Table S4. Structural parameters for 2-Cr and 2-Fe. 
M = Cr or Fe 
2-Cr 
2-Fe 
P–M / Å 
2.3596(8) / 2.3595(7) 
2.3842(8) / 2.4078(7) 
TMPcent···M / Å 
1.739(1) / 1.737(1) 
1.714(1) / 1.711(1) 
TMPcent···M···TMPcent / ° 
177.60(4) 
178.78(2) 
PC2plane···C4planea / ° 
3.57(2) / 4.64(2) 
6.58(2) / 8.73(2) 
P–M–P ring twist / ° 
139.82(5) 
146.35(5) 
a A plane is defined by the three atoms, P and the two adjoining C atoms; a second plane is defined by all four C atoms of the 
TMP ring. The “hinge angle”, HA, between these two planes is reported. 
 
 
P(1)
P(2)
Cr(1)
K(1)
N(1)
N(2)
O(5)
O(6)
O(4)
O(3)
O(1)
O(2)
P(1)
P(2)
Fe(1)
K(1)
N(1)
N(2)
O(5)
O(6)
O(4)
O(3)
O(1)
O(2)


--- Page 13 ---
 
 
S13 
 
 
 
Figure S5. Molecular structure of 1-Cr (left), and the anion of 2-Cr (right) with ellipsoids set at 50%. Hydrogen atoms, methyl 
groups on P(1) ring, and the cation of 2-Cr removed for clarity. Intended to show structural comparisons of the 1,1′-
diphosphametallocene fragment. The pair of numbers for the hinge angle corresponds to the P(1) ring / P(2) ring respectively. 
 
Table S5. Structural comparisons between the 1-Cr and 2-Cr pair. 
M = Cr 
1-Cr 
2-Cr 
Δ 
P–M / Å 
2.3812(4) 
2.3596(8) / 2.3595(7) 
–0.022(1) 
TMPcent···M / Å 
1.795(1) 
1.739(1) / 1.737(1) 
Avg. –0.057(1) 
TMPcent···M···TMPcentb / ° 
180 
177.60(4) 
–2.40° 
PC2plane···C4planea / ° 
4.45(12) 
3.57(2) / 4.64(2) 
c 
a A plane is defined by the three atoms, P and the two adjoining C atoms; a second plane is defined by all four C atoms of the 
TMP ring. The “hinge angle”, HA, between these two planes is reported. 
b The propagated error is too large and precludes meaningful analysis beyond that 2-Cr has approximately the same deviation 
of the P atoms from the TMP plane a 1-Cr. 
 
 
1.795(1) Å
1.795(1) Å
180°
4.45(12)°
P(1)
P(1A)
Cr(1)
P(1)
P(2)
Cr(1)
177.60(4)°
3.57(2)° 
4.64(2)°
1.739(1) Å
1.737(1) Å
1-Cr
2-Cr


--- Page 14 ---
 
 
S14 
 
 
 
Figure S6. Molecular structure of 1-Fe (left) and the anion of 2-Fe (right) with ellipsoids set at 50%. Hydrogen atoms, methyl 
groups on P(1) ring, and the cation of 2-Fe removed for clarity. Intended to show structural comparisons of the 1,1′-
diphosphametallocene fragment. The pair of numbers for the hinge angle corresponds to the P(1) ring / P(2) ring respectively. 
 
Table S6. Structural comparisons between the 1-Fe and 2-Fe pair. 
M = Fe 
1-Fe 
2-Fe 
Δ 
P–M / Å 
2.2932(4) 
2.3842(8) / 2.4078(7) 
+0.103(1) 
TMPcent···M / Å 
1.660(1) 
1.714(1) / 1.711(1) 
Avg. +0.053(1) 
TMPcent···M···TMPcentb / ° 
180 
178.78(2) 
–1.12 
PC2plane···C4planea / ° 
1.27(13) 
6.58(2) / 8.73(2) 
c 
a A plane is defined by the three atoms, P and the two adjoining C atoms; a second plane is defined by all four C atoms of the 
TMP ring. The “hinge angle”, HA, between these two planes is reported. 
b The propagated error is too large and precludes meaningful analysis beyond that 2-Fe has a much larger deviation of the P 
atoms from the TMP plane than 1-Fe. 
 
 
1.660(1) Å
1.660(1) Å
180°
1.27(13)°
P(1)
P(1A)
Fe(1)
P(1)
P(2)
Fe(1)
178.78(2)°
6.58(2)° 
8.73(2)°
1.714(1) Å
1.711(1) Å
1-Fe
2-Fe


--- Page 15 ---
 
 
S15 
 
 
 
Figure S7. Molecular structure of 5 with ellipsoids set at 50%. A disordered component for both Cp rings, and hydrogen atoms 
removed for clarity. Zr(1) sits on a site of symmetry, thus half the molecule is symmetry generated (operations: x, y, z; ½–x, 
y, 1–z). 
Zr(1)···Cpcent = 2.210(5) Å; Zr(1)–C(2) = 2.232(3) Å; C(2)–C(3) = 1.355(4) Å; C(3)–C(3A) = 1.505(6) Å. 
Cpcent···Zr(1) ···Cpcent = 133.6(2)°; C(2)–Zr(1)–C(2A) = 81.22(15)°; Cpcent···Zr(1) ···Cpcent dihedral C(2)–Zr(1)–C(2A) = 
91.66°. 
 
 
Figure S8. Molecular structure of 6 with ellipsoids set at 50%. A second molecule in the asymmetric unit, and hydrogen atoms 
removed for clarity. 
P(1)–C(1) = 1.7987(13) Å; P(1)–C(4) = 1.8057(17) Å; P(1)–C(9) = 1.8313(16) Å; C(1)–C(2) = 1.353(2) Å; C(2)–C(3) = 
1.477(2) Å; C(3)–C(4) = 1.347(2) Å. 
C(1)–P(1)–C(4) = 1.7979(13)°; C(1)–P(1)–C(9) = 106.06(6)°; C(4)–P(1)–C(9) = 103.36(7)°. 
 
 
C(2)
C(2A)
C(3)
C(3A)
Zr(1)
C(4)
C(3)
C(2)
C(1)
P(1)
C(9)


--- Page 16 ---
 
 
S16 
 
 
 
Figure S9. Front view (left) and side view (right) of the molecular structure of 7 with ellipsoids set at 50%, and hydrogen 
atoms removed for clarity. 
Co(1)···Co(1A) = 2.685(1) Å; Co(1)–P(1) = 2.2963(10) Å; Co(1)–P(2) = 2.1697(10) Å; P(1)TMPcent···Co(1) = 1.705(1) Å.  
TPMcent···Co(1···Co(1A) = 166.57(4)°; Co(1)–P(2)–Co(1A) = 76.38(4)°; P(2)–Co(1)–P(2A) = 94.06(4)°; ∑Co(1)–P(2A)–Co(1A)–P(2) 
= 340.88°. 
 
 
Figure S10. Molecular structure of 8 with ellipsoids set at 50%. Hydrogen atoms, a second set of [K(2.2.2-crypt)] and TMP, 
and a THF molecule removed for clarity. The extensive disorder components were also removed, and preclude meaningful 
discussion of the metrical parameters for this compound. 
 
 
P(1)
P(2)
Co(1)
Co(1A)
P(2A)
P(1A)
P(1)
P(2)
Co(1)
Co(1A)
P(2A)
P(1A)
P(1)
N(1)
N(2)
O(1)
O(2)
O(3)
O(4)
K(1)
O(5)
O(6)


--- Page 17 ---
 
 
S17 
 
 
S4. NMR Spectroscopy 
 
Figure S11. 1H NMR spectrum of 1-Fe in C6D6 at room temperature, spectrum cropped to show all observed peaks. * denotes 
a small amount of (TMP)2, presumably formed via oxidative coupling of K-TMP from traces of FeCl3 in commercial samples 
of FeCl2. An authentic sample of (TMP)2 was prepared and the 1H NMR spectra were compared (vide infra, Figure S17) which 
agreed well with literature values.2  
 
 
Figure S12. 31P{1H} NMR spectrum of 1-Fe in C6D6 at room temperature, spectrum cropped to show all observed peaks. 
*denotes a small amount of (TMP)2. 
 
 
1.99, 12 H,
TMP 3,4-Me2
1.64, 12 H, 
d, 3JHP = 9.54 Hz
TMP 2,5-Me2
*
*
Chemical Shift (ppm)
4.0
3.5
3.0
2.5
2.0
1.5
-57.15
*
Chemical Shift (ppm)
40
20
0
-20
-40
-60
-80


--- Page 18 ---
 
 
S18 
 
 
 
Figure S13. 1H NMR spectrum of 2-Cr in C4D8O at room temperature, spectrum cropped to show all observed peaks. 
 
 
Figure S14. 31P{1H} NMR spectrum of 2-Cr in C4D8O at room temperature. The sole observed peak at 74.55 ppm corresponds 
well to the reported 31P resonance for [K(18-crown-6)(TMP)], and is likely a diamagnetic impurity due to decomposition. 
 
 
Chemical Shift (ppm)
8
6
4
2
0
-2
-4
2.63
3.62
3.65
O
O
74.55
Chemical Shift (ppm)
1400
1200
1000
800
600
400
200
0
-200
-400
-600
-800


--- Page 19 ---
 
 
S19 
 
 
 
Figure S15. 1H NMR spectrum of 2-Fe in C4D8O at room temperature, spectrum cropped to show all observed peaks. Broad 
features around –10 ppm, 6 ppm and 12.5 ppm could not be identified by integration or by comparison with known complexes. 
The spectrum can be compared with that of 1-Fe (Figure S11). 
 
 
Figure S16. 31P{1H} NMR spectrum of 2-Fe in C4D8O at room temperature. The sole observed peak at 74.46 ppm corresponds 
well to the reported 31P resonance for [K(18-crown-6)(TMP)], and is likely a diamagnetic impurity due to decomposition. 
 
 
1.94,
TMP 3,4-Me2
2.19, 
d, 3JHP = 10.29 Hz
TMP 2,5-Me2
Chemical Shift (ppm)
14
12
10
8
6
4
2
0
-2
-4
-6
Chemical Shift (ppm)
4.0
3.5
3.0
2.5
2.0
O
O
3.37,
2.2.2-crypt
3.30,
2.2.2-crypt
2.28,
2.2.2-crypt
74.46
Chemical Shift (ppm)
1400
1200
1000
800
600
400
200
0
-200
-400
-600
-800


--- Page 20 ---
 
 
S20 
 
 
 
Figure S17. 1H NMR spectrum of (TMP)2, in C6D6 at room temperature, spectrum cropped to show all observed peaks. The 
peak at 1.89 ppm presents as a triplet but could plausibly be a doublet of doublets due to ABB′ coupling of the CH3 group with 
the 1,3-P and the 1,4-P. 
 
 
Figure S18. 31P{1H} NMR spectrum of (TMP)2, in C6D6 at room temperature. 
 
 
1.73, 12 H,
TMP 3,3',4,4'-Me2
1.89, 12 H, 
m, J = 4.86 / 4.77 Hz
TMP 2,2',5,5'-Me2
Chemical Shift (ppm)
7.0
6.5
6.0
5.5
5.0
4.5
4.0
3.5
3.0
2.5
2.0
-8.43
Chemical Shift (ppm)
140
120
100
80
60
40
20
0
-20
-40
-60
-80


--- Page 21 ---
 
 
S21 
 
 
 
Figure S19. 1H NMR spectrum of 2-Cr in C4D8O, with an external drop of C4H8O outside the FEP liner to act as a reference. 
Collected at room temperature. 
 
 
Figure S20. 1H NMR spectrum of 2-Fe in C4D8O, with an external drop of C4H8O outside the FEP liner to act as a reference. 
Collected at room temperature. The second external THF resonance could not be distinguished from overlapping sample peaks. 
 
Table S7. Solution magnetic susceptibility data for 2-Cr and 2-Fe via Evans NMR method.a 
b Sample / peak 
μeff / B.M 
mol-1 
χ'mT (S.I.)/ 
m3 mol-1 
χ'mT (c.g.s. e.m.u.)/ 
cm3 mol-1 K 
χ'mT (S.I.)/ m3 
mol-1 K 
mass of 
sample / g 
mass of solvent 
+ sample / g 
Mr / g 
mol-1 
Δ peak / 
Hz 
[K(2.2.2-crypt)][Cr(TMP)2] peak 1 
1.455 
1.12E-08 
0.265 
3.33E-06 
0.0150 
0.2465 
745.892 
73.87 
[K(2.2.2-crypt)][Cr(TMP)2] peak 2 
1.433 
1.12E-08 
0.257 
3.33E-06 
0.0150 
0.2465 
745.892 
69.94 
[K(2.2.2-crypt)][Fe(TMP)2] 
1.376 
1.12E-08 
0.237 
3.33E-06 
0.0150 
0.2826 
749.741 
51.77 
a The small masses engender large errors in this methodology, the results should be cautiously interpreted along with other 
data. b Spectrometer frequency 400.130 MHz. Simple diamagnetic correction of Mr / -2,000,000 applied. ρD8THF = 0.985 g mL–
1. 
 
Chemical Shift (ppm)
4.5
4.0
3.5
3.0
2.5
2.0
1.5
1H NMR 400.13 MHz
0.0150 g sample
0.2315 g C4D8O
MW = 745.892 g mol-1
A
B
� = 69.94 Hz
Out
In
O
O
� = 73.87 Hz
Out
In
O
O
1H NMR 400.13 MHz
0.0150 g sample
0.2676 g C4D8O
MW = 749.741 g mol-1
A
B
� = 51.77 Hz
Out
In
O
O
Chemical Shift (ppm)
4.5
4.0
3.5
3.0
2.5
2.0
1.5
1.0


--- Page 22 ---
 
 
S22 
 
 
 
Figure S21. 31P{1H} NMR spectrum of crystals of 7 in C6D6 at room temperature. 
 
 
Chemical Shift (ppm)
220
200
180
160
140
120
100
80
60
40
20
0
96.08
-0.01
-20.37
-40.47
-8.53
(TMP)2


--- Page 23 ---
 
 
S23 
 
 
S5. UV-vis-nIR Spectroscopy 
 
Figure S22. UV-vis-nIR spectrum of 1-Cr in THF (0.63 mM) at room temperature recorded between 5,556–36,364 cm–1 
(1,800–275 nm), window shown between 6,000–32,000 cm–1 (1,667–313 nm). 
 
 
Figure S23. UV-vis-nIR spectrum of 1-Fe in THF (0.50 mM) at room temperature recorded between 5,556–36,364 cm–1 
(1,800–275 nm), window shown between 6,000–32,000 cm–1 (1,667–313 nm). 
 
7,000
11,000
15,000
19,000
23,000
27,000
31,000
1,429
909
667
526
435
370
323
0
1,000
2,000
3,000
4,000
1-Cr
¬ l / nm
e / M-1 cm-1 ®
E / cm-1 ®
16,000
8,000
625
1,250
0
200
400
600
7,000
11,000
15,000
19,000
23,000
27,000
31,000
1,429
909
667
526
435
370
323
0
1,000
2,000
3,000
4,000
5,000
6,000
1-Fe
¬ l / nm
e / M-1 cm-1 ®
E / cm-1 ®
625
1,250
0
200
400
600


--- Page 24 ---
 
 
S24 
 
 
 
Figure S24. UV-vis-nIR spectrum of 2-Cr in THF (0.40 mM) at room temperature recorded between 5,556–36,364 cm–1 
(1,800–275 nm), window shown between 6,000–32,000 cm–1 (1,667–313 nm). Features either side of 7,000 cm–1 are due to 
imperfect background subtraction. 
 
 
Figure S25. UV-vis-nIR spectrum of 2-Fe in THF (0.40 mM) at room temperature recorded between 5,556–36,364 cm–1 
(1,800–275 nm), window shown between 6,000–32,000 cm–1 (1,667–313 nm). Features either side of 7,000 cm–1 are due to 
imperfect background subtraction. 
 
7,000
11,000
15,000
19,000
23,000
27,000
31,000
1,429
909
667
526
435
370
323
0
1,000
2,000
3,000
4,000
5,000
2-Cr
¬ l / nm
e / M-1 cm-1 ®
E / cm-1 ®
16,000
8,000
625
1,250
0
200
400
600
7,000
11,000
15,000
19,000
23,000
27,000
31,000
1,429
909
667
526
435
370
323
0
500
1,000
1,500
2,000
2-Fe
¬ l / nm
e / M-1 cm-1 ®
E / cm-1 ®
16,000
8,000
625
1,250
0
200
400
600


--- Page 25 ---
 
 
S25 
 
 
 
Figure S26. UV-vis-nIR spectrum of both 1-Cr (black, bottom line) and 2-Cr (red, top line) in THF at room temperature 
recorded between 5,556–36,364 cm–1 (1,800–275 nm), window shown between 6,000–32,000 cm–1 (1,667–313 nm). Intended 
to highlight the spectral differences upon reduction. 
 
 
Figure S27. UV-vis-nIR spectrum of both 1-Fe (black, bottom line) and 2-Fe (red, top line) in THF at room temperature 
recorded between 5,556–36,364 cm–1 (1,800–275 nm), window shown between 6,000–32,000 cm–1 (1,667–313 nm). Intended 
to highlight the spectral differences upon reduction. 
 
7,000
11,000
15,000
19,000
23,000
27,000
31,000
1,429
909
667
526
435
370
323
0
1,000
2,000
3,000
4,000
5,000
1-Cr
2-Cr
¬ l / nm
e / M-1 cm-1 ®
E / cm-1 ®
16,000
8,000
625
1,250
0
200
400
600
7,000
11,000
15,000
19,000
23,000
27,000
31,000
1,429
909
667
526
435
370
323
0
1,000
2,000
3,000
4,000
1-Fe
2-Fe
¬ l / nm
e / M-1 cm-1 ®
E / cm-1 ®
625
1,250
200
400
600


--- Page 26 ---
 
 
S26 
 
 
 
Figure S28. A photograph of the solutions used for the UV-vis-nIR measurements. From left to right, 2-Fe, 1-Fe, 2-Cr, and 
1-Cr. 
 
 


--- Page 27 ---
 
 
S27 
 
 
S6. ATR-FTIR Spectroscopy 
 
Figure S29. FT-IR spectrum of both 2-Cr a microcrystalline powder, between 4,000-525 cm–1. 
 
Figure S30. FT-IR spectrum of both 2-Fe a microcrystalline powder, recorded between 4,000-525 cm–1. 
 
 
3,500
3,000
2,500
2,000
1,500
1,000
2-Cr
Transmittance / a.u.
¬ E / cm-1
3,500
3,000
2,500
2,000
1,500
1,000
2-Fe
Transmittance / a.u.
¬ E / cm-1


--- Page 28 ---
 
 
S28 
 
 
 
Figure S31. FT-IR spectrum of both 2-Cr (black, top line) and 2-Fe (red, bottom line) as a microcrystalline powder, recorded 
between 4,000-525 cm–1. Inset shown between 1,600-500 cm–1, to show that the two complexes are essentially identical in the 
fingerprint region. 
 
 
3,500
3,000
2,500
2,000
1,500
1,000
2-Cr
2-Fe
Transmittance / a.u.
¬ E / cm-1
1500
1250
1000
750
Transmittance / a.u.
¬ E / cm-1


--- Page 29 ---
 
 
S29 
 
 
S7. EPR Spectroscopy 
General considerations: High-Field/Frequency Electron Paramagnetic Resonance spectra were recorded on powder samples 
of 2-Cr and 2-Fe. Samples were loaded into polyethylene cups and immobilized with a PTFE stopper. All sample 
manipulations were performed in an inert atmosphere glovebox. The transmission-type spectrometer used in this study 
combined a 17 T superconducting magnet with a phase-locked source (Virginia Diodes Inc., Charlottesville, VA, USA) with 
a series of frequency multipliers. The field modulated signal was detected by an InSb hot-electron bolometer (QMC Ltd., 
Cardiff, U.K.). Temperature control was achieved using an Oxford Instruments (Oxford, U.K.) continuous-flow cryostat. 
Spectral simulations were generated using EasySpin.12 
 
 
Figure S32. Experimental (black) and simulated (red) EPR spectra of 2-Fe (top), 2-Cr (bottom), note the x-axis scale for 2-
Cr. Spectra were recorded at the temperatures and frequencies denoted on the respective figures. Note: the baseline features 
in the 2-Fe spectrum results from micro-crystallites in the imperfectly ground powder. Multi-frequency studies of 2-Cr and 2-
Fe reveal no additional features arising from higher spin states. However, 2-Fe exhibits a minor impurity at lower field which 
is omitted here for clarity and is shown below in Figure S33. The g-values used in the simulations are reported the main text 
as well as in Table S8. 
 
 


--- Page 30 ---
 
 
S30 
 
 
Table S8. Experimental and CASSCF/NEVPT2 Calculated Spin Hamiltonian Parameters. 
 
 
S 
g1, g2, g3 
2-Fe 
Exp. 
1/2 
2.033(5), 1.999(5), 1.943(5) 
NEVPT2 
1/2 
2.116, 2.029,1.919  
2-Cr 
Exp. 
1/2 
2.023(5), 1.993(5), 1.985(5) 
NEVPT2 
1/2 
2.007, 1.984, 1.978 
 
 
Figure S33. Spectrum of 2-Fe (same as Figure S32) with an extended axis to show the unknown minor species. The g-values 
of these features are: ~ 2.09 (*), 2.12 (#), and ~2.18 (†). No evidence of this species was present in 2-Cr. 
 
The sign of Δgμ (gμ − ge), i.e., whether the g – value is less than or greater than the free election value, ge, is determined 
by the competition of the low-energy α- and β- electron transitions. If the excited states contributing to Δg are predominately 
accessed by transitions involving α-electrons, then Δg < 0 while those involving low-lying β-electrons will result in Δg > 0. 
The g-values are reasonably well reproduced by the CASSCF + NEVPT2 calculations (see Table S8) and some qualitative 
insight may be gleaned from examining the AILFT orbital diagrams in Figure 3. In the case of 2-Cr the lowest energy state 
(transitions: {dxy/dy2-y2} ® {dz2}) do not mix by spin-orbit coupling and the sign and magnitude of Δg are determined by the 
mixing of numerous higher energy excited states. This explains the relatively small Δg values (2-Cr, |Δgavg| = 0.0158). For 2-
Fe the first excited state can be described as an α- electron transition from {dxz} ® {dyz} which will mix with the ground state 
and negative Δg along z. As in the case of 2-Cr, several higher energy excited states make up the remaining contributions. 
However, the comparatively large mixing induced by the first excited state is likely a key contributor to the larger observed 
|Δgavg| = 0.0311 in 2-Fe. 
 


--- Page 31 ---
 
 
S31 
 
 
S8. 57Fe Mössbauer Spectroscopy 
General considerations: Spectra were recorded at 120 K in zero applied field using a constant acceleration 
spectrometer and a 57Co/Rh source. The samples used for these measurements consisted of ground powders of 1-
Fe and 2-Fe that were contained in polyethylene sample cups with tightly fitted lids. The isomer shift is reported 
relative to that of α-Fe at room temperature. Spectral simulations were generated using the WMOSS software 
package (SEE Co. Minneapolis, MN).13 
 
Table S9. Experimental and DFT Calculated 57Fe Mössbauer Parameters. 
 
 
d (mm/sec) 
DEQ (mm/sec) 
G (mm/sec)a 
2-Fe 
Exp. 
0.65 
1.28 
-0.602 
DFT 
0.66 
1.46 
- 
1-Fe 
Exp. 
0.48 
2.02 
0.318 
DFT 
0.50 
1.91 
- 
a Positive and negative linewidths indicate Lorentzian and Gaussian lineshapes, respectively. 
 
 


--- Page 32 ---
 
 
S32 
 
 
S9. Quantum Chemical Calculations 
DFT Calculations: Calculations for 1-Cr, 2-Cr, 1-Fe, and 2-Fe were performed using the atomic coordinates determined by 
X-ray crystallography (with hydrogen positions optimized). These calculations were performed using ORCA V4.1 with the 
BP86 functional along with the def2-tzvp basis sets.14-18 The functional was chosen from our previous success in 
reproducing 57Fe Mössbauer parameters using the BP86 functional.19 The resolution of the identity approximation was 
used along with auxiliary basis sets generated using the ‘autoaux’ command.20 The quadrupole splitting and isomer shifts were 
then calculated using the BP86 functional with CP(PPP) (Fe)/def2-tzvp(C)/def2-svp(H). The RI approximation was not used 
for calculation of 57Fe Mössbauer Parameters. The computed density at the Fe nucleus was converted to experimental values 
of the isomer shift using the calibration curve described by Römelt, et al.21 Example input files are shown below. 
 
Complete Active Space Self-Consistent Field/AILFT Calculations: State-averaged complete active space self- consistent 
field (SA-CASSCF) followed by second-order N-electron valence perturbation theory (NEVPT2) calculations were 
performed on the same structures used in the DFT calculations (vide supra). Scalar relativistic effects were 
accounted for by the Douglas−Kroll−Hess (DKH) procedure. SOC was accounted for by quasi-degenerate 
perturbation theory. This methodology of accounting for SOC has proven successful in numerous other studies.22-
30 The dkh-def2-tzvp (Fe,C)/dkh-def2-svp (H) basis set combination was used. The active space consisted of the 
five 3d – orbitals and all of the 3d – electrons (six for Fe2+, seven for Fe+, four for Cr2+, five for Cr+). All roots for 
all multiplicities were included. This means that for Cr2+/Fe2+ 5 quintet, 45 triplet, and 50 singlet states were 
included. For Cr+ 1 sextet, 24 quartet, and 75 singlet states were included. For Fe+ 10 quartet states and 40 doublet 
states were included. The ab-initio ligand field analysis was performed as implemented in ORCA V4.1 and the 
energies presented in the text are those associated with the mapping of the CASSCF+NEVPT2 energies. Example 
input files are shown below.27 
 
 


--- Page 33 ---
 
 
S33 
 
 
S9.1. Example ORCA input files 
Optimize Hydrogen Positions: 
57Fe Mössbauer Parameters: 
! BP86 def2-tzvp autoaux opt 
! BP86 def2-tzvp NoFinalGrid 
%geom 
%method 
optimizehydrogens true 
SpecialGridAtoms 26 
end 
SpecialGridIntAcc 7 
*xyz charge multiplicity 
end 
xyz coordinates of structure 
%basis 
* 
NewGTO H “def2-svp” 
 
NewGTO Fe “CP(PPP)” 
 
end 
 
*xyz charge multiplicity 
 
xyz coordinates of structure 
 
* 
 
%eprnmr 
 
nuclei = all Fe {rho, fgrad} 
 
end 
 
 
 


--- Page 34 ---
 
 
S34 
 
 
CASSCF + NEVPT2: 
!DKH dkh-def2-tzvp autoaux normalprint moread 
%moinp “guess.gbw” # Here the guess orbitals were from a single point DFT calculation 
%casscf # Here the Fe+/d7 settings are used as an example 
nel 7 
norb 5 
mult 4,2 
nroots 10,40 
actorbs dorbs 
trafostep RI 
nevpt2 true 
rel 
dosoc true 
gtensor true 
end 
end 
%basis 
NewGTO H “def2-svp” 
NewGTO Fe “CP(PPP)” 
end 
*xyz -1 2 
xyz coordinates of structure 
* 
 
 
Table S10. Summary of CASSCF+NEVPT2 and AILFT results for 1-Fe, 2-Fe, 1-Cr, and 2-Cr. 
 
Method 
Spin Ground 
State 
AILFT Energies (cm-1) 
1-Fe 
CASSCF 
2 
0.0, 1018.7, 3671.5, 18514.2, 21138.6 
CASSCF+NEVPT
2 
0 
0.0, 899.0, 2984.5, 23425.8, 25707.1 
2-Fe 
CASSCF 
3/2 
0.0, 389.9, 4677.5, 14412.4, 19698.6 
CASSCF+NEVPT
2 
1/2 
0.0, 354.7, 4364.2, 19153.3, 25100.0 
1-Cr 
CASSCF 
2 
0.0, 2226.9, 5170.3, 18790.0, 22484.6 
CASSCF+NEVPT
2 
1 
0.0, 2518.0, 4620.0, 21382.7, 24956.9 
2-Cr 
CASSCF 
1/2 
0.0, 461.3, 7483.5, 21077.6, 25943.8 
CASSCF+NEVPT
2 
1/2 
0.0, 343.7, 6144.3, 24766.6, 29492.1 
 
 


--- Page 35 ---
 
 
S35 
 
 
S10. SQUID Magnetometry 
General considerations: Samples of 1-Cr, 2-Cr, and 2-Fe used for measurements consisted of crushed 
microcrystalline powders weighed (28.0 mg, 1-Cr; 23.6 mg, 2-Cr; 29.0 mg, 2-Fe) into thin-walled precision 
borosilicate NMR tubes, which were flame sealed under vacuum. A diamagnetic correction was performed using 
two further NMR tubes from the same batch also sealed under vacuum. 
Magnetic measurements were performed with a Quantum Design SQUID magnetometer MPMS-XL. The 
temperature dependence of magnetic susceptibility was measured in a direct current (DC) applied magnetic field 
of 0.1 T in the range from 1.8 to 300 K. 
 
Figure S34. Temperature dependence of 𝜒T for 1-Cr and 2-Cr between 1.8 and 300 K. The red lines are to guide 
the eye for values at 298 K (1-Cr, 1.24 emu∙K/mol; 2-Cr, 0.38 emu∙K/mol). 
 
Figure S35. Temperature dependence of 𝜒T for 2-Fe between 1.8 and 300 K. The red line is to guide the eye for 
the value at 298 K (0.26 emu∙K/mol). 
 


--- Page 36 ---
 
 
S36 
 
 
S11. Qualitative ligand HOMOs in Cp– vs {PC4H4}– 
 
Figure S36. Illustrative pictures of the ligand HOMOs relevant to the LUMO in FcH, and [Fe(PC4H4)2]. 
 
The ligand HOMO in Cp– in D5h symmetry is comprised of two degenerate orbitals. However, for {PC4H4}–, the P 
atom breaks the 5-fold symmetry, and thus the analogous orbitals are not degenerate. While both the πP and πC 
orbitals shown above are of appropriate symmetry to interact with the metal dxz, dyz orbitals, they differ in energy 
by ~ 1 eV, and thus the [Fe(PC4H4)2] complex LUMO is mostly comprised of the πP orbital and a metal dyz orbital. 
Please see Fenske et al for a much more in depth analysis,31 as the above discussion does not fully account for the 
lower symmetry of 2-Fe. 
 
 
P
P
e1''
Cp-
�C
�P
Degenerate Cp HOMOs
�P ~ 1 eV lower in energy than �C 
in free {PC4H4}-
{PC4H4}-


--- Page 37 ---
 
 
S37 
 
 
Author Contributions 
CAPG performed all synthesis, UV-vis-nIR and NMR characterization and sample preparation, and X-ray 
diffraction data collection and interpretation with supervision by BLS. SMG performed EPR and Mössbauer 
spectroscopies and interpreted the data, and also performed all theoretical work, with supervision by BWS. OU 
performed SQUID magnetometry measurements and interpretation. RJB sealed samples for analysis and was 
supervised by JLK. CAPG lead the project with project scope and design helped by SMG. The manuscript was 
prepared by CAPG with input from all authors. 
References 
1. 
F. Nief, F. Mathey, L. Ricard, F. Robert, Organometallics, 1988, 7, 921. 
2. 
R. Feher, F. H. Köhler, F. Nief, L. Ricard, S. Rossmayer, Organometallics, 1997, 16, 4606. 
3. 
L. E. Manxzer, J. Deaton, P. Sharp, R. R. Schrock, in Inorganic Syntheses, 2007, pp. 135. 
4. 
T. K. Panda, M. T. Gamer, P. W. Roesky, Organometallics, 2003, 22, 877. 
5. 
P. J. Fagan, W. A. Nugent, J. Am. Chem. Soc., 1988, 110, 2310. 
6. 
P. J. Fagan, W. A. Nugent, Org. Syn., 1992, 70. 
7. 
CrysAlisPro 39.27b, Oxford Diffraction / Agilent Technologies UK Ltd, Yarnton, U.K., 2017. 
8. 
O. V. Dolomanov, L. J. Bourhis, R. J. Gildea, J. A. K. Howard, H. Puschmann, J. Appl. Crystallogr., 2009, 
42, 339. 
9. 
G. M. Sheldrick, Acta Crystallogr. A, 2008, 64, 112. 
10. 
G. M. Sheldrick, Acta Crystallogr. C, 2015, 71, 3. 
11. 
Inkscape: Open Source Scalable Vector Graphics Editor. https://inkscape.org/. 
12. 
S. Stoll, A. Schweiger, J. Magn. Reson., 2006, 178, 42. 
13. 
I. Prisecaru, WMOSS4 Mössbauer Spectral Analysis Software. http://wmoss.org/. 
14. 
A. D. Becke, Phys. Rev. A, 1988, 38, 3098. 
15. 
A. D. Becke, J. Chem. Phys., 1993, 98, 5648. 
16. 
F. Neese, WIRES Comput. Mol. Sci., 2018, 8, e1327. 
17. 
J. P. Perdew, Phys. Rev. B, 1986, 33, 8822. 


--- Page 38 ---
 
 
S38 
 
 
18. 
F. Weigend, R. Ahlrichs, Phys. Chem. Chem. Phys., 2005, 7, 3297. 
19. 
C. A. P. Goodwin, M. Giansiracusa, S. M. Greer, H. M. Nicholas, P. Evans, M. Vonci, S. Hill, N. F. 
Chilton, D. P. Mills, Nat. Chem., 2020. 
20. 
G. L. Stoychev, A. A. Auer, F. Neese, J. Chem. Theory. Comput., 2017, 13, 554. 
21. 
M. Römelt, S. Ye, F. Neese, Inorg. Chem., 2009, 48, 784. 
22. 
S. K. Singh, J. Eng, M. Atanasov, F. Neese, Coord. Chem. Rev., 2017, 344, 2. 
23. 
D. Maganas, J. Krzystek, E. Ferentinos, A. M. Whyte, N. Robertson, V. Psycharis, A. Terzis, F. Neese, P. 
Kyritsis, Inorg. Chem., 2012, 51, 7218. 
24. 
S. Ye, F. Neese, J. Chem. Theory. Comput., 2012, 8, 2344. 
25. 
J. Jung, M. Atanasov, F. Neese, Inorg. Chem., 2017, 56, 8802. 
26. 
E. A. Suturina, J. Nehrkorn, J. M. Zadrozny, J. Liu, M. Atanasov, T. Weyhermuller, D. Maganas, S. Hill, 
A. Schnegg, E. Bill, J. R. Long, F. Neese, Inorg. Chem., 2017, 56, 3102. 
27. 
M. Atanasov, D. Aravena, E. Suturina, E. Bill, D. Maganas, F. Neese, Coord. Chem. Rev., 2015, 289-290, 
177. 
28. 
E. A. Suturina, D. Maganas, E. Bill, M. Atanasov, F. Neese, Inorg. Chem., 2015, 54, 9948. 
29. 
M. Autillo, M. A. Islam, J. Jung, J. Pilme, N. Galland, L. Guerin, P. Moisy, C. Berthon, C. Tamain, H. 
Bolvin, Phys. Chem. Chem. Phys., 2020, 22, 14293. 
30. 
J. Jung, M. A. Islam, V. L. Pecoraro, T. Mallah, C. Berthon, H. Bolvin, Chemistry, 2019, 25, 15112. 
31. 
N. M. Kostic, R. F. Fenske, Organometallics, 1983, 2, 1008. 
 
 


--- Page 39 ---
 
 
S39 
 
 
S13. CIF Reports 
 
checkCIF/PLATON report 
Structure factors have been supplied for datablock(s) 1-Cr
THIS REPORT IS FOR GUIDANCE ONLY. IF USED AS PART OF A REVIEW PROCEDURE
FOR PUBLICATION, IT SHOULD NOT REPLACE THE EXPERTISE OF AN EXPERIENCED
CRYSTALLOGRAPHIC REFEREE.
No syntax errors found.        CIF dictionary        Interpreting this report
Datablock: 1-Cr 
Bond precision:
C-C = 0.0021 A
Wavelength=0.71073
Cell:
a=7.8420(3)
b=12.5620(4)
c=8.8611(3)
alpha=90
beta=109.444(4)
gamma=90
Temperature: 
100 K
Calculated
Reported
Volume
823.13(5)
823.13(5)
Space group
P 21/n 
P 1 21/n 1 
Hall group
-P 2yn 
-P 2yn 
Moiety formula
C16 H24 Cr P2
C16 H24 Cr P2
Sum formula
C16 H24 Cr P2
C16 H24 Cr P2
Mr
330.29
330.29
Dx,g cm-3
1.333
1.333
Z
2
2
Mu (mm-1)
0.874
0.874
F000
348.0
348.0
F000’
349.12
h,k,lmax
9,15,10
9,15,10
Nref
1507 
1508 
Tmin,Tmax
0.700,0.769
0.804,1.000
Tmin’
0.686
Correction method= # Reported T Limits: Tmin=0.804 Tmax=1.000
AbsCorr = MULTI-SCAN
Data completeness= 1.001
Theta(max)= 25.349
R(reflections)= 0.0252( 1441)
wR2(reflections)= 0.0667( 1508)
S = 1.091
Npar= 92
The following ALERTS were generated. Each ALERT has the format
       test-name_ALERT_alert-type_alert-level.
Click on the hyperlinks for more details of the test.


--- Page 40 ---
 
 
S40 
 
 
 
 Alert level G
PLAT328_ALERT_4_G Possible Missing H on sp3? Phosphorus ..........         P1 Check 
PLAT380_ALERT_4_G Incorrectly? Oriented X(sp2)-Methyl Moiety .....         C6 Check 
PLAT883_ALERT_1_G No Info/Value for _atom_sites_solution_primary .     Please Do !  
PLAT941_ALERT_3_G Average HKL Measurement Multiplicity ...........        4.5 Low   
PLAT978_ALERT_2_G Number C-C Bonds with Positive Residual Density.          6 Info  
   0 ALERT level A = Most likely a serious problem - resolve or explain
   0 ALERT level B = A potentially serious problem, consider carefully
   0 ALERT level C = Check. Ensure it is not caused by an omission or oversight
   5 ALERT level G = General information/check it is not something unexpected
   1 ALERT type 1 CIF construction/syntax error, inconsistent or missing data
   1 ALERT type 2 Indicator that the structure model may be wrong or deficient
   1 ALERT type 3 Indicator that the structure quality may be low
   2 ALERT type 4 Improvement, methodology, query or suggestion
   0 ALERT type 5 Informative message, check
It is advisable to attempt to resolve as many as possible of the alerts in all categories. Often the
minor alerts point to easily fixed oversights, errors and omissions in your CIF or refinement
strategy, so attention to these fine details can be worthwhile. In order to resolve some of the more
serious problems it may be necessary to carry out additional measurements or structure
refinements. However, the purpose of your study may justify the reported deviations and the more
serious of these should normally be commented upon in the discussion or experimental section of a
paper or in the "special_details" fields of the CIF. checkCIF was carefully designed to identify
outliers and unusual parameters, but every test has its limitations and alerts that are not important
in a particular case may appear. Conversely, the absence of alerts does not guarantee there are no
aspects of the results needing attention. It is up to the individual to critically assess their own
results and, if necessary, seek expert advice.
Publication of your CIF in IUCr journals 
A basic structural check has been run on your CIF. These basic checks will be run on all CIFs
submitted for publication in IUCr journals (Acta Crystallographica, Journal of Applied 
Crystallography, Journal of Synchrotron Radiation); however, if you intend to submit to Acta
Crystallographica Section C or E or IUCrData, you should make sure that full publication checks
are run on the final version of your CIF prior to submission.
Publication of your CIF in other journals 
Please refer to the Notes for Authors of the relevant journal for any special instructions relating to
CIF submission.
PLATON version of 16/07/2020; check.def file version of 12/07/2020 


--- Page 41 ---
 
 
S41 
 
 
 
checkCIF/PLATON report 
Structure factors have been supplied for datablock(s) 1-Fe
THIS REPORT IS FOR GUIDANCE ONLY. IF USED AS PART OF A REVIEW PROCEDURE
FOR PUBLICATION, IT SHOULD NOT REPLACE THE EXPERTISE OF AN EXPERIENCED
CRYSTALLOGRAPHIC REFEREE.
No syntax errors found.        CIF dictionary        Interpreting this report
Datablock: 1-Fe 
Bond precision:
C-C = 0.0030 A
Wavelength=0.71073
Cell:
a=14.2948(7)
b=12.8549(6)
c=8.8448(4)
alpha=90
beta=104.364(5)
gamma=90
Temperature: 
100 K
Calculated
Reported
Volume
1574.50(13)
1574.50(13)
Space group
C 2/c 
C 1 2/c 1 
Hall group
-C 2yc 
-C 2yc 
Moiety formula
C16 H24 Fe P2
C16 H24 Fe P2
Sum formula
C16 H24 Fe P2
C16 H24 Fe P2
Mr
334.14
334.14
Dx,g cm-3
1.410
1.410
Z
4
4
Mu (mm-1)
1.146
1.146
F000
704.0
704.0
F000’
706.33
h,k,lmax
17,16,11
17,16,11
Nref
1608 
1605 
Tmin,Tmax
0.770,0.832
0.850,1.000
Tmin’
0.717
Correction method= # Reported T Limits: Tmin=0.850 Tmax=1.000
AbsCorr = MULTI-SCAN
Data completeness= 0.998
Theta(max)= 26.370
R(reflections)= 0.0312( 1407)
wR2(reflections)= 0.0798( 1605)
S = 1.117
Npar= 92
The following ALERTS were generated. Each ALERT has the format
       test-name_ALERT_alert-type_alert-level.
Click on the hyperlinks for more details of the test.


--- Page 42 ---
 
 
S42 
 
 
 
 Alert level C
PLAT094_ALERT_2_C Ratio of Maximum / Minimum Residual Density ....       2.29 Report
PLAT911_ALERT_3_C Missing FCF Refl Between Thmin & STh/L=    0.600          4 Report
 Alert level G
PLAT328_ALERT_4_G Possible Missing H on sp3? Phosphorus ..........         P1 Check 
PLAT941_ALERT_3_G Average HKL Measurement Multiplicity ...........        3.9 Low   
PLAT978_ALERT_2_G Number C-C Bonds with Positive Residual Density.          6 Info  
   0 ALERT level A = Most likely a serious problem - resolve or explain
   0 ALERT level B = A potentially serious problem, consider carefully
   2 ALERT level C = Check. Ensure it is not caused by an omission or oversight
   3 ALERT level G = General information/check it is not something unexpected
   0 ALERT type 1 CIF construction/syntax error, inconsistent or missing data
   2 ALERT type 2 Indicator that the structure model may be wrong or deficient
   2 ALERT type 3 Indicator that the structure quality may be low
   1 ALERT type 4 Improvement, methodology, query or suggestion
   0 ALERT type 5 Informative message, check
It is advisable to attempt to resolve as many as possible of the alerts in all categories. Often the
minor alerts point to easily fixed oversights, errors and omissions in your CIF or refinement
strategy, so attention to these fine details can be worthwhile. In order to resolve some of the more
serious problems it may be necessary to carry out additional measurements or structure
refinements. However, the purpose of your study may justify the reported deviations and the more
serious of these should normally be commented upon in the discussion or experimental section of a
paper or in the "special_details" fields of the CIF. checkCIF was carefully designed to identify
outliers and unusual parameters, but every test has its limitations and alerts that are not important
in a particular case may appear. Conversely, the absence of alerts does not guarantee there are no
aspects of the results needing attention. It is up to the individual to critically assess their own
results and, if necessary, seek expert advice.
Publication of your CIF in IUCr journals 
A basic structural check has been run on your CIF. These basic checks will be run on all CIFs
submitted for publication in IUCr journals (Acta Crystallographica, Journal of Applied 
Crystallography, Journal of Synchrotron Radiation); however, if you intend to submit to Acta
Crystallographica Section C or E or IUCrData, you should make sure that full publication checks
are run on the final version of your CIF prior to submission.
Publication of your CIF in other journals 
Please refer to the Notes for Authors of the relevant journal for any special instructions relating to
CIF submission.
PLATON version of 16/07/2020; check.def file version of 12/07/2020 


--- Page 43 ---
 
 
S43 
 
 
 
checkCIF/PLATON report 
Structure factors have been supplied for datablock(s) 2-Cr
THIS REPORT IS FOR GUIDANCE ONLY. IF USED AS PART OF A REVIEW PROCEDURE
FOR PUBLICATION, IT SHOULD NOT REPLACE THE EXPERTISE OF AN EXPERIENCED
CRYSTALLOGRAPHIC REFEREE.
No syntax errors found.        CIF dictionary        Interpreting this report
Datablock: 2-Cr 
Bond precision:
C-C = 0.0030 A
Wavelength=0.71073
Cell:
a=10.1994(4)
b=14.2193(7)
c=15.0882(6)
alpha=62.328(4)
beta=86.597(3)
gamma=87.156(3)
Temperature: 
100 K
Calculated
Reported
Volume
1933.91(16)
1933.91(16)
Space group
P -1 
P -1 
Hall group
-P 1 
-P 1 
Moiety formula
C18 H36 K N2 O6, C16 H24
Cr P2
C16 H24 Cr P2, C18 H36 K
N2 O6
Sum formula
C34 H60 Cr K N2 O6 P2
C34 H60 Cr K N2 O6 P2
Mr
745.88
745.88
Dx,g cm-3
1.281
1.281
Z
2
2
Mu (mm-1)
0.529
0.529
F000
798.0
798.0
F000’
799.73
h,k,lmax
12,17,18
12,17,18
Nref
7083 
7060 
Tmin,Tmax
0.853,0.853
0.902,1.000
Tmin’
0.853
Correction method= # Reported T Limits: Tmin=0.902 Tmax=1.000
AbsCorr = MULTI-SCAN
Data completeness= 0.997
Theta(max)= 25.350
R(reflections)= 0.0338( 5848)
wR2(reflections)= 0.0872( 7060)
S = 1.100
Npar= 423


--- Page 44 ---
 
 
S44 
 
 
 
The following ALERTS were generated. Each ALERT has the format
       test-name_ALERT_alert-type_alert-level.
Click on the hyperlinks for more details of the test.
 Alert level C
PLAT911_ALERT_3_C Missing FCF Refl Between Thmin & STh/L=    0.600         13 Report
 Alert level G
PLAT042_ALERT_1_G Calc. and Reported MoietyFormula Strings  Differ     Please Check 
PLAT328_ALERT_4_G Possible Missing H on sp3? Phosphorus ..........         P1 Check 
PLAT328_ALERT_4_G Possible Missing H on sp3? Phosphorus ..........         P2 Check 
PLAT380_ALERT_4_G Incorrectly? Oriented X(sp2)-Methyl Moiety .....         C7 Check 
PLAT380_ALERT_4_G Incorrectly? Oriented X(sp2)-Methyl Moiety .....        C16 Check 
PLAT910_ALERT_3_G Missing # of FCF Reflection(s) Below Theta(Min).          1 Note  
PLAT912_ALERT_4_G Missing # of FCF Reflections Above STh/L=  0.600         10 Note  
PLAT933_ALERT_2_G Number of OMIT Records in Embedded .res File ...          1 Note  
PLAT941_ALERT_3_G Average HKL Measurement Multiplicity ...........        3.3 Low   
PLAT978_ALERT_2_G Number C-C Bonds with Positive Residual Density.          8 Info  
   0 ALERT level A = Most likely a serious problem - resolve or explain
   0 ALERT level B = A potentially serious problem, consider carefully
   1 ALERT level C = Check. Ensure it is not caused by an omission or oversight
  10 ALERT level G = General information/check it is not something unexpected
   1 ALERT type 1 CIF construction/syntax error, inconsistent or missing data
   2 ALERT type 2 Indicator that the structure model may be wrong or deficient
   3 ALERT type 3 Indicator that the structure quality may be low
   5 ALERT type 4 Improvement, methodology, query or suggestion
   0 ALERT type 5 Informative message, check


--- Page 45 ---
 
 
S45 
 
 
 
It is advisable to attempt to resolve as many as possible of the alerts in all categories. Often the
minor alerts point to easily fixed oversights, errors and omissions in your CIF or refinement
strategy, so attention to these fine details can be worthwhile. In order to resolve some of the more
serious problems it may be necessary to carry out additional measurements or structure
refinements. However, the purpose of your study may justify the reported deviations and the more
serious of these should normally be commented upon in the discussion or experimental section of a
paper or in the "special_details" fields of the CIF. checkCIF was carefully designed to identify
outliers and unusual parameters, but every test has its limitations and alerts that are not important
in a particular case may appear. Conversely, the absence of alerts does not guarantee there are no
aspects of the results needing attention. It is up to the individual to critically assess their own
results and, if necessary, seek expert advice.
Publication of your CIF in IUCr journals 
A basic structural check has been run on your CIF. These basic checks will be run on all CIFs
submitted for publication in IUCr journals (Acta Crystallographica, Journal of Applied 
Crystallography, Journal of Synchrotron Radiation); however, if you intend to submit to Acta
Crystallographica Section C or E or IUCrData, you should make sure that full publication checks
are run on the final version of your CIF prior to submission.
Publication of your CIF in other journals 
Please refer to the Notes for Authors of the relevant journal for any special instructions relating to
CIF submission.
PLATON version of 16/07/2020; check.def file version of 12/07/2020 


--- Page 46 ---
 
 
S46 
 
 
 
checkCIF/PLATON report 
Structure factors have been supplied for datablock(s) 2-Fe
THIS REPORT IS FOR GUIDANCE ONLY. IF USED AS PART OF A REVIEW PROCEDURE
FOR PUBLICATION, IT SHOULD NOT REPLACE THE EXPERTISE OF AN EXPERIENCED
CRYSTALLOGRAPHIC REFEREE.
No syntax errors found.        CIF dictionary        Interpreting this report
Datablock: 2-Fe 
Bond precision:
C-C = 0.0034 A
Wavelength=0.71073
Cell:
a=10.1947(2)
b=14.2373(4)
c=15.0677(3)
alpha=62.275(2)
beta=85.840(2)
gamma=87.045(2)
Temperature: 
100 K
Calculated
Reported
Volume
1930.42(8)
1930.42(8)
Space group
P -1 
P -1 
Hall group
-P 1 
-P 1 
Moiety formula
C18 H36 K N2 O6, C16 H24
Fe P2
C16 H24 Fe P2, C18 H36 K
N2 O6
Sum formula
C34 H60 Fe K N2 O6 P2
C34 H60 Fe K N2 O6 P2
Mr
749.73
749.73
Dx,g cm-3
1.290
1.290
Z
2
2
Mu (mm-1)
0.624
0.624
F000
802.0
802.0
F000’
803.77
h,k,lmax
12,17,18
12,17,18
Nref
7890 
7807 
Tmin,Tmax
0.856,0.969
0.872,1.000
Tmin’
0.856
Correction method= # Reported T Limits: Tmin=0.872 Tmax=1.000
AbsCorr = MULTI-SCAN
Data completeness= 0.989
Theta(max)= 26.370
R(reflections)= 0.0400( 6351)
wR2(reflections)= 0.1035( 7807)
S = 1.090
Npar= 423


--- Page 47 ---
 
 
S47 
 
 
 
The following ALERTS were generated. Each ALERT has the format
       test-name_ALERT_alert-type_alert-level.
Click on the hyperlinks for more details of the test.
 Alert level C
PLAT094_ALERT_2_C Ratio of Maximum / Minimum Residual Density ....       2.41 Report
PLAT906_ALERT_3_C Large K Value in the Analysis of Variance ......      2.110 Check 
PLAT911_ALERT_3_C Missing FCF Refl Between Thmin & STh/L=    0.600         60 Report
 Alert level G
PLAT042_ALERT_1_G Calc. and Reported MoietyFormula Strings  Differ     Please Check 
PLAT154_ALERT_1_G The s.u.’s on the Cell Angles are Equal ..(Note)      0.002 Degree
PLAT328_ALERT_4_G Possible Missing H on sp3? Phosphorus ..........         P1 Check 
PLAT328_ALERT_4_G Possible Missing H on sp3? Phosphorus ..........         P2 Check 
PLAT883_ALERT_1_G No Info/Value for _atom_sites_solution_primary .     Please Do !  
PLAT910_ALERT_3_G Missing # of FCF Reflection(s) Below Theta(Min).          1 Note  
PLAT912_ALERT_4_G Missing # of FCF Reflections Above STh/L=  0.600         22 Note  
PLAT933_ALERT_2_G Number of OMIT Records in Embedded .res File ...          1 Note  
PLAT941_ALERT_3_G Average HKL Measurement Multiplicity ...........        3.3 Low   
PLAT978_ALERT_2_G Number C-C Bonds with Positive Residual Density.          6 Info  
   0 ALERT level A = Most likely a serious problem - resolve or explain
   0 ALERT level B = A potentially serious problem, consider carefully
   3 ALERT level C = Check. Ensure it is not caused by an omission or oversight
  10 ALERT level G = General information/check it is not something unexpected
   3 ALERT type 1 CIF construction/syntax error, inconsistent or missing data
   3 ALERT type 2 Indicator that the structure model may be wrong or deficient
   4 ALERT type 3 Indicator that the structure quality may be low
   3 ALERT type 4 Improvement, methodology, query or suggestion
   0 ALERT type 5 Informative message, check


--- Page 48 ---
 
 
S48 
 
 
 
It is advisable to attempt to resolve as many as possible of the alerts in all categories. Often the
minor alerts point to easily fixed oversights, errors and omissions in your CIF or refinement
strategy, so attention to these fine details can be worthwhile. In order to resolve some of the more
serious problems it may be necessary to carry out additional measurements or structure
refinements. However, the purpose of your study may justify the reported deviations and the more
serious of these should normally be commented upon in the discussion or experimental section of a
paper or in the "special_details" fields of the CIF. checkCIF was carefully designed to identify
outliers and unusual parameters, but every test has its limitations and alerts that are not important
in a particular case may appear. Conversely, the absence of alerts does not guarantee there are no
aspects of the results needing attention. It is up to the individual to critically assess their own
results and, if necessary, seek expert advice.
Publication of your CIF in IUCr journals 
A basic structural check has been run on your CIF. These basic checks will be run on all CIFs
submitted for publication in IUCr journals (Acta Crystallographica, Journal of Applied 
Crystallography, Journal of Synchrotron Radiation); however, if you intend to submit to Acta
Crystallographica Section C or E or IUCrData, you should make sure that full publication checks
are run on the final version of your CIF prior to submission.
Publication of your CIF in other journals 
Please refer to the Notes for Authors of the relevant journal for any special instructions relating to
CIF submission.
PLATON version of 16/07/2020; check.def file version of 12/07/2020 


--- Page 49 ---
 
 
S49 
 
 
 
checkCIF/PLATON report 
Structure factors have been supplied for datablock(s) 5
THIS REPORT IS FOR GUIDANCE ONLY. IF USED AS PART OF A REVIEW PROCEDURE
FOR PUBLICATION, IT SHOULD NOT REPLACE THE EXPERTISE OF AN EXPERIENCED
CRYSTALLOGRAPHIC REFEREE.
No syntax errors found.        CIF dictionary        Interpreting this report
Datablock: 5 
Bond precision:
C-C = 0.0040 A
Wavelength=0.71073
Cell:
a=14.1554(8)
b=9.6976(4)
c=12.0182(6)
alpha=90
beta=112.336(6)
gamma=90
Temperature: 
150 K
Calculated
Reported
Volume
1526.00(15)
1525.99(14)
Space group
I 2/a 
I 1 2/a 1 
Hall group
-I 2ya 
-I 2ya 
Moiety formula
C18 H22 Zr
C18 H22 Zr
Sum formula
C18 H22 Zr
C18 H22 Zr
Mr
329.58
329.57
Dx,g cm-3
1.434
1.435
Z
4
4
Mu (mm-1)
0.705
0.705
F000
680.0
680.0
F000’
668.18
h,k,lmax
17,12,15
17,12,15
Nref
1559 
1558 
Tmin,Tmax
0.905,0.974
0.775,1.000
Tmin’
0.889
Correction method= # Reported T Limits: Tmin=0.775 Tmax=1.000
AbsCorr = GAUSSIAN
Data completeness= 0.999
Theta(max)= 26.386
R(reflections)= 0.0330( 1401)
wR2(reflections)= 0.0805( 1558)
S = 1.068
Npar= 117
The following ALERTS were generated. Each ALERT has the format
       test-name_ALERT_alert-type_alert-level.
Click on the hyperlinks for more details of the test.


--- Page 50 ---
 
 
S50 
 
 
 
 Alert level G
PLAT003_ALERT_2_G Number of Uiso or Uij Restrained non-H Atoms ...         15 Report
PLAT187_ALERT_4_G The CIF-Embedded .res File Contains RIGU Records          1 Report
PLAT301_ALERT_3_G Main Residue  Disorder ..............(Resd  1  )        53% Note  
PLAT811_ALERT_5_G No ADDSYM Analysis: Too Many Excluded Atoms ....          ! Info  
PLAT860_ALERT_3_G Number of Least-Squares Restraints .............        144 Note  
PLAT883_ALERT_1_G No Info/Value for _atom_sites_solution_primary .     Please Do !  
PLAT910_ALERT_3_G Missing # of FCF Reflection(s) Below Theta(Min).          1 Note  
PLAT912_ALERT_4_G Missing # of FCF Reflections Above STh/L=  0.600          1 Note  
PLAT941_ALERT_3_G Average HKL Measurement Multiplicity ...........        1.8 Low   
PLAT955_ALERT_1_G Reported (CIF) and Actual (FCF) Lmax Differ by .          1 Units 
PLAT978_ALERT_2_G Number C-C Bonds with Positive Residual Density.          8 Info  
   0 ALERT level A = Most likely a serious problem - resolve or explain
   0 ALERT level B = A potentially serious problem, consider carefully
   0 ALERT level C = Check. Ensure it is not caused by an omission or oversight
  11 ALERT level G = General information/check it is not something unexpected
   2 ALERT type 1 CIF construction/syntax error, inconsistent or missing data
   2 ALERT type 2 Indicator that the structure model may be wrong or deficient
   4 ALERT type 3 Indicator that the structure quality may be low
   2 ALERT type 4 Improvement, methodology, query or suggestion
   1 ALERT type 5 Informative message, check
It is advisable to attempt to resolve as many as possible of the alerts in all categories. Often the
minor alerts point to easily fixed oversights, errors and omissions in your CIF or refinement
strategy, so attention to these fine details can be worthwhile. In order to resolve some of the more
serious problems it may be necessary to carry out additional measurements or structure
refinements. However, the purpose of your study may justify the reported deviations and the more
serious of these should normally be commented upon in the discussion or experimental section of a
paper or in the "special_details" fields of the CIF. checkCIF was carefully designed to identify
outliers and unusual parameters, but every test has its limitations and alerts that are not important
in a particular case may appear. Conversely, the absence of alerts does not guarantee there are no
aspects of the results needing attention. It is up to the individual to critically assess their own
results and, if necessary, seek expert advice.
Publication of your CIF in IUCr journals 
A basic structural check has been run on your CIF. These basic checks will be run on all CIFs
submitted for publication in IUCr journals (Acta Crystallographica, Journal of Applied 
Crystallography, Journal of Synchrotron Radiation); however, if you intend to submit to Acta
Crystallographica Section C or E or IUCrData, you should make sure that full publication checks
are run on the final version of your CIF prior to submission.
Publication of your CIF in other journals 
Please refer to the Notes for Authors of the relevant journal for any special instructions relating to
CIF submission.
PLATON version of 16/07/2020; check.def file version of 12/07/2020 


--- Page 51 ---
 
 
S51 
 
 
 
checkCIF/PLATON report 
Structure factors have been supplied for datablock(s) 6
THIS REPORT IS FOR GUIDANCE ONLY. IF USED AS PART OF A REVIEW PROCEDURE
FOR PUBLICATION, IT SHOULD NOT REPLACE THE EXPERTISE OF AN EXPERIENCED
CRYSTALLOGRAPHIC REFEREE.
No syntax errors found.        CIF dictionary        Interpreting this report
Datablock: 6 
Bond precision:
C-C = 0.0020 A
Wavelength=1.54184
Cell:
a=8.0727(2)
b=8.74422(19)
c=18.4065(4)
alpha=103.1739(17)
beta=92.6248(19)
gamma=100.705(2)
Temperature: 150 K
Calculated
Reported
Volume
1237.81(5)
1237.81(5)
Space group
P -1 
P -1 
Hall group
-P 1 
-P 1 
Moiety formula
C14 H17 P
C14 H17 P
Sum formula
C14 H17 P
C14 H17 P
Mr
216.25
216.24
Dx,g cm-3
1.160
1.160
Z
4
4
Mu (mm-1)
1.664
1.664
F000
464.0
464.0
F000’
466.15
h,k,lmax
9,10,22
9,10,22
Nref
4541 
4482 
Tmin,Tmax
0.760,0.862
0.703,1.000
Tmin’
0.690
Correction method= # Reported T Limits: Tmin=0.703 Tmax=1.000
AbsCorr = GAUSSIAN
Data completeness= 0.987
Theta(max)= 68.247
R(reflections)= 0.0329( 4211)
wR2(reflections)= 0.0899( 4482)
S = 1.036
Npar= 279
The following ALERTS were generated. Each ALERT has the format
       test-name_ALERT_alert-type_alert-level.
Click on the hyperlinks for more details of the test.


--- Page 52 ---
 
 
S52 
 
 
 
 Alert level C
PLAT911_ALERT_3_C Missing FCF Refl Between Thmin & STh/L=    0.600         55 Report
 Alert level G
PLAT412_ALERT_2_G Short Intra XH3 .. XHn     H6A      ..H5D      .       2.02 Ang.  
                                                      x,y,z  =      1_555 Check 
PLAT912_ALERT_4_G Missing # of FCF Reflections Above STh/L=  0.600          5 Note  
PLAT941_ALERT_3_G Average HKL Measurement Multiplicity ...........        2.9 Low   
PLAT978_ALERT_2_G Number C-C Bonds with Positive Residual Density.          8 Info  
PLAT992_ALERT_5_G Repd & Actual _reflns_number_gt Values Differ by          2 Check 
   0 ALERT level A = Most likely a serious problem - resolve or explain
   0 ALERT level B = A potentially serious problem, consider carefully
   1 ALERT level C = Check. Ensure it is not caused by an omission or oversight
   5 ALERT level G = General information/check it is not something unexpected
   0 ALERT type 1 CIF construction/syntax error, inconsistent or missing data
   2 ALERT type 2 Indicator that the structure model may be wrong or deficient
   2 ALERT type 3 Indicator that the structure quality may be low
   1 ALERT type 4 Improvement, methodology, query or suggestion
   1 ALERT type 5 Informative message, check
It is advisable to attempt to resolve as many as possible of the alerts in all categories. Often the
minor alerts point to easily fixed oversights, errors and omissions in your CIF or refinement
strategy, so attention to these fine details can be worthwhile. In order to resolve some of the more
serious problems it may be necessary to carry out additional measurements or structure
refinements. However, the purpose of your study may justify the reported deviations and the more
serious of these should normally be commented upon in the discussion or experimental section of a
paper or in the "special_details" fields of the CIF. checkCIF was carefully designed to identify
outliers and unusual parameters, but every test has its limitations and alerts that are not important
in a particular case may appear. Conversely, the absence of alerts does not guarantee there are no
aspects of the results needing attention. It is up to the individual to critically assess their own
results and, if necessary, seek expert advice.
Publication of your CIF in IUCr journals 
A basic structural check has been run on your CIF. These basic checks will be run on all CIFs
submitted for publication in IUCr journals (Acta Crystallographica, Journal of Applied 
Crystallography, Journal of Synchrotron Radiation); however, if you intend to submit to Acta
Crystallographica Section C or E or IUCrData, you should make sure that full publication checks
are run on the final version of your CIF prior to submission.
Publication of your CIF in other journals 
Please refer to the Notes for Authors of the relevant journal for any special instructions relating to
CIF submission.
PLATON version of 16/07/2020; check.def file version of 12/07/2020 


--- Page 53 ---
 
 
S53 
 
 
 
checkCIF/PLATON report 
Structure factors have been supplied for datablock(s) 7
THIS REPORT IS FOR GUIDANCE ONLY. IF USED AS PART OF A REVIEW PROCEDURE
FOR PUBLICATION, IT SHOULD NOT REPLACE THE EXPERTISE OF AN EXPERIENCED
CRYSTALLOGRAPHIC REFEREE.
No syntax errors found.        CIF dictionary        Interpreting this report
Datablock: 7 
Bond precision:
C-C = 0.0040 A
Wavelength=0.71073
Cell:
a=20.161(3)
b=9.0620(4)
c=20.214(3)
alpha=90
beta=119.71(2)
gamma=90
Temperature: 
100 K
Calculated
Reported
Volume
3207.6(9)
3207.7(9)
Space group
C 2/c 
C 1 2/c 1 
Hall group
-C 2yc 
-C 2yc 
Moiety formula
C32 H48 Co2 P4
C32 H48 Co2 P4
Sum formula
C32 H48 Co2 P4
C32 H48 Co2 P4
Mr
674.44
674.44
Dx,g cm-3
1.397
1.397
Z
4
4
Mu (mm-1)
1.254
1.254
F000
1416.0
1416.0
F000’
1420.68
h,k,lmax
24,10,24
24,10,24
Nref
2938 
2935 
Tmin,Tmax
0.835,0.882
0.861,1.000
Tmin’
0.606
Correction method= # Reported T Limits: Tmin=0.861 Tmax=1.000
AbsCorr = MULTI-SCAN
Data completeness= 0.999
Theta(max)= 25.342
R(reflections)= 0.0333( 2325)
wR2(reflections)= 0.0834( 2935)
S = 1.009
Npar= 180
The following ALERTS were generated. Each ALERT has the format
       test-name_ALERT_alert-type_alert-level.
Click on the hyperlinks for more details of the test.


--- Page 54 ---
 
 
S54 
 
 
 
 Alert level G
PLAT232_ALERT_2_G Hirshfeld Test Diff (M-X)  Co1      --P1       .        5.5 s.u.  
PLAT328_ALERT_4_G Possible Missing H on sp3? Phosphorus ..........         P1 Check 
PLAT883_ALERT_1_G No Info/Value for _atom_sites_solution_primary .     Please Do !  
PLAT910_ALERT_3_G Missing # of FCF Reflection(s) Below Theta(Min).          1 Note  
PLAT912_ALERT_4_G Missing # of FCF Reflections Above STh/L=  0.600          1 Note  
PLAT933_ALERT_2_G Number of OMIT Records in Embedded .res File ...          1 Note  
PLAT978_ALERT_2_G Number C-C Bonds with Positive Residual Density.          4 Info  
   0 ALERT level A = Most likely a serious problem - resolve or explain
   0 ALERT level B = A potentially serious problem, consider carefully
   0 ALERT level C = Check. Ensure it is not caused by an omission or oversight
   7 ALERT level G = General information/check it is not something unexpected
   1 ALERT type 1 CIF construction/syntax error, inconsistent or missing data
   3 ALERT type 2 Indicator that the structure model may be wrong or deficient
   1 ALERT type 3 Indicator that the structure quality may be low
   2 ALERT type 4 Improvement, methodology, query or suggestion
   0 ALERT type 5 Informative message, check
It is advisable to attempt to resolve as many as possible of the alerts in all categories. Often the
minor alerts point to easily fixed oversights, errors and omissions in your CIF or refinement
strategy, so attention to these fine details can be worthwhile. In order to resolve some of the more
serious problems it may be necessary to carry out additional measurements or structure
refinements. However, the purpose of your study may justify the reported deviations and the more
serious of these should normally be commented upon in the discussion or experimental section of a
paper or in the "special_details" fields of the CIF. checkCIF was carefully designed to identify
outliers and unusual parameters, but every test has its limitations and alerts that are not important
in a particular case may appear. Conversely, the absence of alerts does not guarantee there are no
aspects of the results needing attention. It is up to the individual to critically assess their own
results and, if necessary, seek expert advice.
Publication of your CIF in IUCr journals 
A basic structural check has been run on your CIF. These basic checks will be run on all CIFs
submitted for publication in IUCr journals (Acta Crystallographica, Journal of Applied 
Crystallography, Journal of Synchrotron Radiation); however, if you intend to submit to Acta
Crystallographica Section C or E or IUCrData, you should make sure that full publication checks
are run on the final version of your CIF prior to submission.
Publication of your CIF in other journals 
Please refer to the Notes for Authors of the relevant journal for any special instructions relating to
CIF submission.
PLATON version of 16/07/2020; check.def file version of 12/07/2020 


--- Page 55 ---
 
 
S55 
 
 
 
checkCIF/PLATON report 
Structure factors have been supplied for datablock(s) 8
THIS REPORT IS FOR GUIDANCE ONLY. IF USED AS PART OF A REVIEW PROCEDURE
FOR PUBLICATION, IT SHOULD NOT REPLACE THE EXPERTISE OF AN EXPERIENCED
CRYSTALLOGRAPHIC REFEREE.
No syntax errors found.        CIF dictionary        Interpreting this report
Datablock: 8 
Bond precision:
C-C = 0.0050 A
Wavelength=0.71073
Cell:
a=12.7189(8)
b=15.9437(10)
c=17.8937(12)
alpha=64.066(6)
beta=85.655(5)
gamma=86.321(5)
Temperature: 
100 K
Calculated
Reported
Volume
3251.9(4)
3251.9(4)
Space group
P -1 
P -1 
Hall group
-P 1 
-P 1 
Moiety formula
2(C18 H36 K N2 O6), 2(C8
H12 P), C4 H8 O
2(C8 H12 P), C4 H8 O,
2(C18 H36 K N2 O6)
Sum formula
C56 H104 K2 N4 O13 P2
C56 H104 K2 N4 O13 P2
Mr
1181.57
1181.57
Dx,g cm-3
1.207
1.207
Z
2
2
Mu (mm-1)
0.254
0.254
F000
1280.0
1280.0
F000’
1281.74
h,k,lmax
15,19,22
15,19,22
Nref
13295 
13040 
Tmin,Tmax
0.958,0.987
0.826,1.000
Tmin’
0.910
Correction method= # Reported T Limits: Tmin=0.826 Tmax=1.000
AbsCorr = MULTI-SCAN
Data completeness= 0.981
Theta(max)= 26.372
R(reflections)= 0.0673( 8417)
wR2(reflections)= 0.1979( 13040)
S = 1.063
Npar= 1310


--- Page 56 ---
 
 
S56 
 
 
 
The following ALERTS were generated. Each ALERT has the format
       test-name_ALERT_alert-type_alert-level.
Click on the hyperlinks for more details of the test.
 Alert level C
PLAT088_ALERT_3_C Poor Data / Parameter Ratio ....................       9.95 Note  
PLAT329_ALERT_4_C Carbon Atom Hybridisation Unclear for ..........       C35B Check 
PLAT340_ALERT_3_C Low Bond Precision on  C-C Bonds ...............      0.005 Ang.  
PLAT906_ALERT_3_C Large K Value in the Analysis of Variance ......      3.712 Check 
PLAT911_ALERT_3_C Missing FCF Refl Between Thmin & STh/L=    0.600        168 Report
PLAT977_ALERT_2_C Check Negative Difference Density on H6B              -0.35 eA-3  
PLAT977_ALERT_2_C Check Negative Difference Density on H6C              -0.40 eA-3  
 Alert level G
PLAT002_ALERT_2_G Number of Distance or Angle Restraints on AtSite         10 Note  
PLAT003_ALERT_2_G Number of Uiso or Uij Restrained non-H Atoms ...        136 Report
PLAT042_ALERT_1_G Calc. and Reported MoietyFormula Strings  Differ     Please Check 
PLAT176_ALERT_4_G The CIF-Embedded .res File Contains SADI Records          3 Report
PLAT178_ALERT_4_G The CIF-Embedded .res File Contains SIMU Records          5 Report
PLAT301_ALERT_3_G Main Residue  Disorder ..............(Resd  1  )       100% Note  
PLAT301_ALERT_3_G Main Residue  Disorder ..............(Resd  3  )       100% Note  
PLAT301_ALERT_3_G Main Residue  Disorder ..............(Resd  4  )       100% Note  
PLAT302_ALERT_4_G Anion/Solvent/Minor-Residue Disorder (Resd  2  )       100% Note  
PLAT302_ALERT_4_G Anion/Solvent/Minor-Residue Disorder (Resd  6  )       100% Note  
PLAT302_ALERT_4_G Anion/Solvent/Minor-Residue Disorder (Resd  7  )       100% Note  
PLAT302_ALERT_4_G Anion/Solvent/Minor-Residue Disorder (Resd  8  )       100% Note  
PLAT302_ALERT_4_G Anion/Solvent/Minor-Residue Disorder (Resd  9  )       100% Note  
PLAT304_ALERT_4_G Non-Integer Number of Atoms in ..... (Resd  1  )      57.33 Check 
PLAT304_ALERT_4_G Non-Integer Number of Atoms in ..... (Resd  2  )      55.50 Check 
PLAT304_ALERT_4_G Non-Integer Number of Atoms in ..... (Resd  3  )       5.67 Check 
PLAT304_ALERT_4_G Non-Integer Number of Atoms in ..... (Resd  4  )       7.50 Check 
PLAT304_ALERT_4_G Non-Integer Number of Atoms in ..... (Resd  6  )      17.14 Check 
PLAT304_ALERT_4_G Non-Integer Number of Atoms in ..... (Resd  7  )       3.86 Check 
PLAT304_ALERT_4_G Non-Integer Number of Atoms in ..... (Resd  8  )      10.85 Check 
PLAT304_ALERT_4_G Non-Integer Number of Atoms in ..... (Resd  9  )       2.14 Check 
PLAT380_ALERT_4_G Incorrectly? Oriented X(sp2)-Methyl Moiety .....         C5 Check 
PLAT380_ALERT_4_G Incorrectly? Oriented X(sp2)-Methyl Moiety .....         C6 Check 
PLAT380_ALERT_4_G Incorrectly? Oriented X(sp2)-Methyl Moiety .....         C7 Check 
PLAT380_ALERT_4_G Incorrectly? Oriented X(sp2)-Methyl Moiety .....         C8 Check 
PLAT380_ALERT_4_G Incorrectly? Oriented X(sp2)-Methyl Moiety .....       C31A Check 
PLAT380_ALERT_4_G Incorrectly? Oriented X(sp2)-Methyl Moiety .....       C32A Check 
PLAT380_ALERT_4_G Incorrectly? Oriented X(sp2)-Methyl Moiety .....       C33A Check 
PLAT380_ALERT_4_G Incorrectly? Oriented X(sp2)-Methyl Moiety .....       C34A Check 
PLAT398_ALERT_2_G Deviating  C-O-C    Angle From 120 for O13A           100.9 Degree
PLAT398_ALERT_2_G Deviating  C-O-C    Angle From 120 for O13B            87.6 Degree
PLAT413_ALERT_2_G Short Inter XH3 .. XHn     H5A      ..H53C     .       2.09 Ang.  
                                                1-x,1-y,1-z  =      2_666 Check 
PLAT413_ALERT_2_G Short Inter XH3 .. XHn     H6C      ..H42A     .       2.09 Ang.  
                                                 1-x,-y,1-z  =      2_656 Check 
PLAT413_ALERT_2_G Short Inter XH3 .. XHn     H6C      ..H42C     .       2.09 Ang.  
                                                 1-x,-y,1-z  =      2_656 Check 
PLAT720_ALERT_4_G Number of Unusual/Non-Standard Labels ..........          4 Note  
PLAT773_ALERT_2_G Check long C-C Bond in CIF: C35B     --C36B            1.78 Ang.  
PLAT773_ALERT_2_G Check long C-C Bond in CIF: C45B     --C46B            1.86 Ang.  
PLAT790_ALERT_4_G Centre of Gravity not Within Unit Cell: Resd.  #          2 Note  
              C18 H36 K N2 O6                                                   
PLAT790_ALERT_4_G Centre of Gravity not Within Unit Cell: Resd.  #          4 Note  
              C18 H36 K N2 O6                                                   
PLAT811_ALERT_5_G No ADDSYM Analysis: Too Many Excluded Atoms ....          ! Info  


--- Page 57 ---
 
 
S57 
 
 
 
 
PLAT860_ALERT_3_G Number of Least-Squares Restraints .............       2121 Note  
PLAT910_ALERT_3_G Missing # of FCF Reflection(s) Below Theta(Min).          4 Note  
PLAT912_ALERT_4_G Missing # of FCF Reflections Above STh/L=  0.600         84 Note  
PLAT933_ALERT_2_G Number of OMIT Records in Embedded .res File ...          4 Note  
PLAT941_ALERT_3_G Average HKL Measurement Multiplicity ...........        2.4 Low   
PLAT978_ALERT_2_G Number C-C Bonds with Positive Residual Density.          5 Info  
PLAT992_ALERT_5_G Repd & Actual _reflns_number_gt Values Differ by          1 Check 
   0 ALERT level A = Most likely a serious problem - resolve or explain
   0 ALERT level B = A potentially serious problem, consider carefully
   7 ALERT level C = Check. Ensure it is not caused by an omission or oversight
  47 ALERT level G = General information/check it is not something unexpected
   1 ALERT type 1 CIF construction/syntax error, inconsistent or missing data
  13 ALERT type 2 Indicator that the structure model may be wrong or deficient
  10 ALERT type 3 Indicator that the structure quality may be low
  28 ALERT type 4 Improvement, methodology, query or suggestion
   2 ALERT type 5 Informative message, check
It is advisable to attempt to resolve as many as possible of the alerts in all categories. Often the
minor alerts point to easily fixed oversights, errors and omissions in your CIF or refinement
strategy, so attention to these fine details can be worthwhile. In order to resolve some of the more
serious problems it may be necessary to carry out additional measurements or structure
refinements. However, the purpose of your study may justify the reported deviations and the more
serious of these should normally be commented upon in the discussion or experimental section of a
paper or in the "special_details" fields of the CIF. checkCIF was carefully designed to identify
outliers and unusual parameters, but every test has its limitations and alerts that are not important
in a particular case may appear. Conversely, the absence of alerts does not guarantee there are no
aspects of the results needing attention. It is up to the individual to critically assess their own
results and, if necessary, seek expert advice.
Publication of your CIF in IUCr journals 
A basic structural check has been run on your CIF. These basic checks will be run on all CIFs
submitted for publication in IUCr journals (Acta Crystallographica, Journal of Applied 
Crystallography, Journal of Synchrotron Radiation); however, if you intend to submit to Acta
Crystallographica Section C or E or IUCrData, you should make sure that full publication checks
are run on the final version of your CIF prior to submission.
Publication of your CIF in other journals 
Please refer to the Notes for Authors of the relevant journal for any special instructions relating to
CIF submission.
PLATON version of 16/07/2020; check.def file version of 12/07/2020